\documentclass[12pt]{book}

\usepackage[T1]{fontenc}
\usepackage[utf8]{inputenc}
\usepackage{microtype}
\usepackage{setspace}
\usepackage{geometry}
\usepackage{longtable}
\usepackage{lipsum}
\usepackage{csquotes}
\usepackage{hyperref}

\geometry{margin=1in}
\setstretch{1.15}

% Long s substitution (ſ)
% (We will typeset ſ directly in the text.)

\begin{document}

%%%%%%%%%%%%%%%%%%%%%%%%%%%%%%%%%%%%%%%%%%%%%%%%%%%%%%%%%%%%
%%%%%%%%%%%%%%%%%%%%%%%% TITLE PAGE %%%%%%%%%%%%%%%%%%%%%%%%%
%%%%%%%%%%%%%%%%%%%%%%%%%%%%%%%%%%%%%%%%%%%%%%%%%%%%%%%%%%%%

\thispagestyle{empty}

\begin{center}

{\Large\bfseries THE KANT–SWEDENBORG BOUNDARY PAPERS}\[1em]
{\large A CRITICAL EDITION OF CERTAIN UNPUBLISHED MANUSCRIPTS,\
JOURNALS, AND PRIVATE LETTERS}\[2em]

{\large Collected and Tranſcribed from Authentic Sources}\[2em]

{\large \itshape Printed for the Learned Reader}\[1em]

{\large \itshape By G.~E.~Hartmann, Printer in Leipzig}\[2em]

{\large MDCCXC}

\end{center}

\newpage
\thispagestyle{empty}

\begin{center}
\itshape
``Reaſon is not a Window but a Boundary.''\[1em]
\upshape
--- A marginal Note found amongſt the Papers of Immanuel Kant (Nachlaſs~K–42)\[2em]

\itshape
``The Lattice trembleth, and the World with it.''\[1em]
\upshape
--- From the laſt Journal of Emanuel Swedenborg, dated March~1772
\end{center}

\newpage

%%%%%%%%%%%%%%%%%%%%%%%%%%%%%%%%%%%%%%%%%%%%%%%%%%%%%%%%%%%%
%%%%%%%%%%%%%%%%%%%%%%%% DEDICATION %%%%%%%%%%%%%%%%%%%%%%%%%
%%%%%%%%%%%%%%%%%%%%%%%%%%%%%%%%%%%%%%%%%%%%%%%%%%%%%%%%%%%%

\chapter*{Dedication}

To the ſtudents and lovers of natural and moral philoſophy, and to all who hold dear the labour of the enquiring Mind, this Volume is humbly dedicated; not that it preſumes to teach, but that it may provoke thought, and perhaps excite in the ſerious Reader a tremor of that ſacred Curioſity which ſtands at the brink of Knowledge, yet dares not paſs beyond the appointed Boundary.

%%%%%%%%%%%%%%%%%%%%%%%%%%%%%%%%%%%%%%%%%%%%%%%%%%%%%%%%%%%%
%%%%%%%%%%%%%%%%%%%%%%%% PREFACE %%%%%%%%%%%%%%%%%%%%%%%%%
%%%%%%%%%%%%%%%%%%%%%%%%%%%%%%%%%%%%%%%%%%%%%%%%%%%%%%%%%%%%

\chapter*{Preface of the Editor}

The Editor, in offering the preſent Collection of Papers to the Publick, is not inſenſible of the hazard he incurreth. For the Materials herein contained partake of a moſt uncommon Character, being neither mere trifles of Curioſity nor yet the plain and orderly Concluſions of eſtabliſhed Science. Rather they conſiſt of private Notes, loſt Journals, and fragments of Correſpondence which paſſed betwixt two Men of extraordinary Genius: Immanuel~Kant and Emanuel~Swedenborg.

The tranſcription of theſe Papers hath not been undertaken lightly. From the moment their Exiſtence firſt became known to the Editor, he was haunted by a Doubt whether they ought ever to be publiſhed. For they reveal, or ſeem to reveal, Occurrences of a moſt peculiar and potent Nature during the Years~1768 to~1772, at which Time both Kant and Swedenborg underwent Changes ſo marked and ſo profound that their reſpective Writings before and after thoſe Years may almoſt be regarded as Works of two different Minds.

It cannot eſcape the attenſion of any careful Reader that the Year~1769 ſtands as a kind of hinge or turning; a Year of inward Diſquiet, of trembling Thought, of viſionary or perhaps delirious Reflection. Though the Editor doth not preſume to determine the Cauſe of theſe inward Alterations, he deemeth it his Duty to bring theſe Documents to Light, that the Publick may judge for itſelf.

And if the Reader find among theſe Leaves expreſſions which appear dark, obſcure, or even terrifying, let it be remembered that they were not compoſed with a View to Publication. They are the private utterances of Minds ſtrained to their extremity, labouring under Conceptions which they could neither wholly comprehend nor wholly reject.

The Editor hath aimed to preſerve the orthography and punctuation of the Originals, except where abſolute Neceſſity demanded a gentle Correction. In all other matters he hath preferred Fidelity to Elegance.

May Providence guide the Reader’s Judgement.

\hfill ---,The Editor\
\hfill Leipzig, Michaelmas Term, MDCCXC

%%%%%%%%%%%%%%%%%%%%%%%%%%%%%%%%%%%%%%%%%%%%%%%%%%%%%%%%%%%%
%%%%%%%%%%%%%%%%%%%%%%% CONTENTS %%%%%%%%%%%%%%%%%%%%%%%%%%%
%%%%%%%%%%%%%%%%%%%%%%%%%%%%%%%%%%%%%%%%%%%%%%%%%%%%%%%%%%%%

\tableofcontents

%%%%%%%%%%%%%%%%%%%%%%%%%%%%%%%%%%%%%%%%%%%%%%%%%%%%%%%%%%%%
%%%%%%%%%%%%%%%%%%%%%%%% PART I %%%%%%%%%%%%%%%%%%%%%%%%%%%%%
%%%%%%%%%%%%%%%%%%%%%%%%%%%%%%%%%%%%%%%%%%%%%%%%%%%%%%%%%%%%

\part*{Part~I\
A Hiſtorical Introduction to the Years~1768–1772}

\addcontentsline{toc}{part}{Part~I: A Hiſtorical Introduction to the Years~1768–1772}

%%%%%%%%%%%%%%%%%%%%%%%%%%%%%%%%%%%%%%%%%%%%%%%%%%%%%%%%%%%%
%%%%%%%%%%%%%%%%%%%%%%% CHAPTER I %%%%%%%%%%%%%%%%%%%%%%%%%%
%%%%%%%%%%%%%%%%%%%%%%%%%%%%%%%%%%%%%%%%%%%%%%%%%%%%%%%%%%%%

\chapter*{Chapter~I\[0.5em]
Concerning the Condition of Philoſophical Enquiry in the Late Decades of the Century}

\addcontentsline{toc}{chapter}{Chapter~I: Concerning the Condition of Philoſophical Enquiry}

In the latter Decades of this enlightened Century, Philoſophy found itſelf in a ſtrange Twilight; for though the Lamp of Reaſon burned with a moſt conſiderable Brightneſs, yet the Shadings around that Light were deep and unexamined. Many laboured under the Conviction that the Bounds of human Knowledge might eaſily be advanced by the mere Extension of known Methods; while others, more cautious, ſuſpected that there exiſted certain Limits beyond which the Mind might not paſs without incurring Diſorder.

Amidſt this Scene roſe the Figure of Immanuel~Kant, who in his earlier Writings had embraced the received Metaphyſicks of his Age, but who, in the Years immediately preceding his Inaugural Diſſertation, began to perceive that the Foundations of all Knowledge reſted not upon the Nature of Things as they are in themſelves, but upon the inward Structures by which the Mind apprehendeth them. His Thought, in this Period, may be likened to a Sea growing troubled, not by outward Storms, but by ſubterranean Motions hidden from all Eyes but his own.

It was likewiſe in this Period that the aged Swedenborg, already renowned for his viſionary Writings and for his curious Inſights into the Nature of the Spirit, began to record in his private Journals certain Experiences of a Character uncommonly vivid and weighty, wherein the Boundaries between the ſenſible and the ſupersenſible ſeemed to tremble.

It is the Conjunction of theſe two Figures, and the inward Agitations they endured in the ſame narrow Span of Years, that inviteth our moſt diligent Enquiry.

%%%%%%%%%%%%%%%%%%%%%%%%%%%%%%%%%%%%%%%%%%%%%%%%%%%%%%%%%%%%
%%%%%%%%%%%%%%%%%%%%%%% CHAPTER II %%%%%%%%%%%%%%%%%%%%%%%%%
%%%%%%%%%%%%%%%%%%%%%%%%%%%%%%%%%%%%%%%%%%%%%%%%%%%%%%%%%%%%

\chapter*{Chapter~II\[0.5em]
On the Lives and Circumſtances of Kant and Swedenborg before the Year~1769}

\addcontentsline{toc}{chapter}{Chapter~II: On the Lives and Circumſtances}

Immanuel~Kant, born in the City of Königsberg in the Year~1724, paſſed his early Life in Study and Modeſt Labour, devoting himſelf to the Sciences of Nature and the mathematical Arts, thereby forming that exact and penetrating Habit of Mind for which he afterwards became ſo juſtly celebrated. His Lectures were much frequented, and his Diſcourſe was ſingular in its Clarity and quiet Force.

Emanuel~Swedenborg, elder by ſeveral Decades, had begun his Career as a Man of natural Science, and in his early Treatiſes had diſplayed a Genius no leſs acute than Kant's, though directed to other Ends. In the latter Part of his Life, however, he turned his Attention to viſionary Subjects, affirming that he had been favoured with Inſights into the hidden Realms of Spirit.

Before the Year~1769, there is no Record of any direct Intercourſe between theſe two remarkable Men. Yet the Seeds of their later Conjunction may already be traced in their Writings: in Swedenborg’s growing Preoccupation with the Structures underlying viſible Reality, and in Kant’s mounting Suſpicion that the Forms of Intuition and the Categories of the Underſtanding were themſelves the Conditions by which the World of Appearance is upheld.

%%%%%%%%%%%%%%%%%%%%%%%%%%%%%%%%%%%%%%%%%%%%%%%%%%%%%%%%%%%%
%%%%%%%%%%%%%%%%%%%%%%% CHAPTER III %%%%%%%%%%%%%%%%%%%%%%%%%
%%%%%%%%%%%%%%%%%%%%%%%%%%%%%%%%%%%%%%%%%%%%%%%%%%%%%%%%%%%%

\chapter*{Chapter~III\[0.5em]
On the Notable Winter of the Year~1769}

\addcontentsline{toc}{chapter}{Chapter~III: On the Notable Winter of 1769}

The Winter of the Year~1769 hath, in the Teſtimonies here collected, a moſt peculiar Significance. For in that Seaſon there appear, both in Kant’s private Notes and in Swedenborg’s Journal, Marks of a diſordered and fearful Experience encountered by each Man in his own Way. Certain Paſſages ſpeak of a trembling in the inward Structures of Perception; others of Viſions, or Viſitations, which neither of them could wholly comprehend, nor yet deny.

It was during this Winter that Kant began to withdraw from his cuſtomary Lectures, making Excuses in his Letters which the Editor, upon Inſpection, findeth uncharacteristically vague. Swedenborg, for his Part, recordeth Viſions of a City inverted beneath a River’s Surface, of Shadows that moved contrary to their Objects, and of a Voice—or perhaps a Structure—which uttered the dread Injunction: Write the Boundary.

That theſe Experiences occurred in the ſame Months is a Fact which cannot eaſily be laid aſide.

%%%%%%%%%%%%%%%%%%%%%%%%%%%%%%%%%%%%%%%%%%%%%%%%%%%%%%%%%%%%
%%%%%%%%%%%%%%%%%%%%%%% CHAPTER IV %%%%%%%%%%%%%%%%%%%%%%%%%
%%%%%%%%%%%%%%%%%%%%%%%%%%%%%%%%%%%%%%%%%%%%%%%%%%%%%%%%%%%%

\chapter*{Chapter~IV\[0.5em]
On the Fragmentary Nature of the Surviving Evidence}

\addcontentsline{toc}{chapter}{Chapter~IV: On the Fragmentary Nature of the Evidence}

It is the misfortune of the Editor that the Papers here tranſcribed have not come down to us in a perfect State. Many Leaves of Kant’s Nachlaſs were found worn, torn, or ſcorched by Fire; ſeveral of Swedenborg’s Entries were blotted, as if by Tears or by the trembling Hand of an aged Writer. Some Letters bear no Date or Signature, and muſt be placed with Care and Conjecture.

Yet the very Fragmentari­neſs of theſe Documents lendeth them an Air of inward Truth; for they appear not as the contrived Fictions of later Times, but as the ſpontaneous Outpourings of Minds encountering a Burden too heavy to be borne in the publick Eye.

In compiling this Volume, the Editor hath taken every Precaution to preſerve what remained, ſupplying no more than was abſolutely required for the Sake of plain Underſtanding. Where the Text ſeemeth obſcure, it hath been left ſo. Where it trembleth, the Tremor hath not been ſmoothed.

For whatſoever did shake theſe two great Minds, in that Winter ſo many Years ago, hath left its Mark as surely in the broken Sentences as in the more perfect Paragraphs.

\newpage

\part*{Part~II[0.5em]
Metaphyſical and Critical Reflections upon the Nature of Appearance and the Hidden Ground}

\addcontentsline{toc}{part}{Part~II: Metaphyſical and Critical Reflections}

%%%%%%%%%%%%%%%%%%%%%%%%%%%%%%%%%%%%%%%%%%%%%%%%%%%%%%%%%%%%
%%%%%%%%%%%%%%%%%%%%%%%% CHAPTER V %%%%%%%%%%%%%%%%%%%%%%%%%
%%%%%%%%%%%%%%%%%%%%%%%%%%%%%%%%%%%%%%%%%%%%%%%%%%%%%%%%%%%%

\chapter*{Chapter~V[0.5em]
Concerning the Ancient Queſtion of the Thing–in–Itſelf}

\addcontentsline{toc}{chapter}{Chapter~V: Concerning the Thing–in–Itſelf}

It hath long been known to the Students of Philoſophy that the Mind, in its moſt ſtrict Enquiries, diſtinguiſheth between the Object as it appears and the Object as it may be conceived to exiſt in its own hidden Eſſence. This Diſtinction, familiar to the Ancient Sages and renewed in later Centuries, became in the Decade here examined a Matter of urgent concern; for the inward Obſervations recorded by the two principal Figures of this Volume ſeem to preſs, with unwonted Weight, upon the Border betwixt Appearance and that more terrible Ground which underlieth it.

Immanuel~Kant, even in the early Fragments tranſcribed from this Period, employeth a Manner of Speech moſt uncharacteriſtic of his maturer Writings. He ſpeaketh of a \textit{veil}, of a \textit{boundary}, of a \textit{frame and lattice} by which all Appearance is upheld, and of the lamentable Conſequence that would enſue if this Lattice were ever to tremble or be removed. For a Man of his diſciplined Temper, trained in the moſt exact Sciences, ſuch Language ſtriketh the Reader as thoſe firſt dark glimmerings of a Thought that hath not yet found its Order.

Swedenborg, meanwhile, affirmeth in his Journals of the ſame Months that he beheld fleeting Viſions of certain Structures anterior to the World of Senſe; and though the Editor maketh no Claim to determine the Nature of theſe Viſions, yet it is moſt remarkable that the inward Forms Swedenborg laboureth to deſcribe bear a curious Reſemblance to Kant’s own later Doctrine of the Categories and of the Forms of Intuition.

Thus, even before the Critique had taken its final Shape, we behold the firſt Faintneſs of a Thought which, like an infant Creature, ſtruggled to be born.

%%%%%%%%%%%%%%%%%%%%%%%%%%%%%%%%%%%%%%%%%%%%%%%%%%%%%%%%%%%%
%%%%%%%%%%%%%%%%%%%%%%%% CHAPTER VI %%%%%%%%%%%%%%%%%%%%%%%%%
%%%%%%%%%%%%%%%%%%%%%%%%%%%%%%%%%%%%%%%%%%%%%%%%%%%%%%%%%%%%

\chapter*{Chapter~VI[0.5em]
On the Precarious Foundation of Space and Time}

\addcontentsline{toc}{chapter}{Chapter~VI: On the Foundation of Space and Time}

It would be an Error to imagine that the Conceptions of Space and Time, which we handle with ſuch familiar Readineſs, poſſeſs a Stability independent of the Mind’s own Conſtitution. For in the Papers of the Year~1769, we encounter the ſingular Notion—firſt hinted in Kant’s private Notes, and confirmed by Swedenborg’s viſionary Leaves—that Space and Time are not infinite Oceans into which the Mind is thrown, but delicate Surfaces drawn taut over a deeper Abyſs.

Certain Paſſages from the Winter of that Year bear Witneſs to unuſual Experiences: the perceived Stoppage of a Town Clock for the Space of a Breath, the Duplication of an Object’s Fall, or the Inverſion of familiar Streets beneath the Surface of a frozen River. Though theſe Reports are doubtleſs coloured by the inward Diſpoſition of the Writers, their Coincidence admoniſheth us that the Structures by which we apprehend the World are more frail than common Opinion ſuppoſeth.

In one Fragment Kant writeth that `Time is a Ribbon laid upon the Mind, and if the Hand that draweth it were to falter, all Succeſſion would overlap.'' Another Note, evidently from the ſame Night, declareth: `Space is a Lattice; its Lines may be ſeen when the Eye trembleth.'' Such Expreſſions, though obſcure, reveal an inward Agitation, as though the Man had drawn nearer than he wiſhed to the Border where Perception turneth upon its own Foundations.

Swedenborg, for his Part, ſeemed to apprehend the ſame Truth from the oppoſite Direction, affirming that `the Hall of Appearance ſtandeth upon Beams not visible to the Eye of Fleſh,'' and that theſe Beams `shift and incline when the Mind is weary or when the World itſelf is in Diſquiet.'' Whatever may be thought of his Viſions, they contain a Strain of Conception not wholly alien to the Critick who later proclaimed that Space and Time are Forms of the Intuition, and not Attributes of Things as they are in themſelves.

Thus, the more we contemplate the teſtimonies of theſe Troubled Years, the more we are drawn toward the Concluſion that the Notions of Space and Time, far from being impregnable, are maintained by a delicate Harmony of the Mind’s own Force—a Harmony which, in rare and dreadful Moments, may fall into Diſcord.

%%%%%%%%%%%%%%%%%%%%%%%%%%%%%%%%%%%%%%%%%%%%%%%%%%%%%%%%%%%%
%%%%%%%%%%%%%%%%%%%%%%%% CHAPTER VII %%%%%%%%%%%%%%%%%%%%%%%%
%%%%%%%%%%%%%%%%%%%%%%%%%%%%%%%%%%%%%%%%%%%%%%%%%%%%%%%%%%%%

\chapter*{Chapter~VII[0.5em]
On the Diſturbance of the Categories and the Dread Conſequences Thereof}

\addcontentsline{toc}{chapter}{Chapter~VII: On the Diſturbance of the Categories}

Among the Papers herein collected, none are more ſtriking than thoſe which ſpeak of a certain Diſturbance or Torsion in the inward Categories by which the Mind arrangeth the Manifold of Senſation. For in his later Writings, Kant defineth the Categories as pure Conceptions of the Underſtanding, fixed and immutable, and the very Conditions of any poſſible Experience. Yet in his private Notes from the Year~1769, we perceive a moſt alarming Conceit: namely, that the Categories themſelves may be ſubject to alteration, fluctuation, or trembling.

One Margin containeth this fearful Obſervation: `If Cauſality were to waver, even for the ſpace of a Moment, the World would ſpilt like heated Metal and ſwallow its own Order.'' In another, diſtinguiſhed by a Hand more urgent and leſs controlled: `The Shadows moved before the Objects; therefore Succession is not ſecure.''

Contemporaneous Paſſages in Swedenborg’s Journals exhibit a like Anxiety. He recordeth Viſions of `moving Forms anterior to Objects,'' of `Chains that turned before the Links,'' of ``Appearances that ſtood while their Shadows walked.'' Though one may interpret theſe Deſcriptions as the Productions of an overwrought Fancy, yet their Conformity with Kant’s inward Terrours ſeemeth too exact to be coincidental.

Whatſoever their Cauſe, theſe Teſtimonies admoniſh us that the inward Architecture of Experience—ſo firm in the Daylight of common Perception—may, under certain unknown Conditions, diſplay a frailty terrifying to the Mind that beholdeth it.

And it is perhaps no Wonder that Kant, having ſurveilled this Abyſs, ſhould thereafter labour with a moſt inexorable Severity to erect a System of Critick wherein the very Notion of cauſal or categorical Diſorder is rendered impious, if not inconceivable. For it is the Fate of many a profound Thinker, that he ſeeketh to bind with Iron the very Gates that once opened to him in a Moment of inward Diſturbance.

%%%%%%%%%%%%%%%%%%%%%%%%%%%%%%%%%%%%%%%%%%%%%%%%%%%%%%%%%%%%
%%%%%%%%%%%%%%%%%%%%%%%% CHAPTER VIII %%%%%%%%%%%%%%%%%%%%%%%
%%%%%%%%%%%%%%%%%%%%%%%%%%%%%%%%%%%%%%%%%%%%%%%%%%%%%%%%%%%%

\chapter*{Chapter~VIII[0.5em]
On the Convergence of Two Divergent Minds}

\addcontentsline{toc}{chapter}{Chapter~VIII: On the Convergence of Two Minds}

Nothing in the Paſt Lives of either Man prepared the World for the brief but momentous Convergence that occurred between Kant and Swedenborg in the Years here examined. Each was of a Nature diametrically oppoſed to the other: Kant, the ſevere Critick of Reaſon, who ſought to limit the Pretences of the Mind in order to ſecure its moſt certain Grounds; Swedenborg, the viſionary Seer, who affirmed that the Mind might pierce the moſt inviolate Veil of Nature.

Yet in theſe Years of inward Turmoil, the two Minds drew near upon a Ground neither of them had choſen. For Kant, in the fragmentary Notes tranſcribed here, expreſſeth Thoughts of a moſt uncommon Nature, wherein the edgy Border between Phenomenon and Noumenon ſeemed to waver. Swedenborg, in his final Journals, declareth that he had perceived `another Labourer at the Boundary,'' who `ſtood with trembling Hand, fearing to draw the final Line.''

It is not for the Editor to determine the full Meaning of theſe Paſſages; but they appear to record a momentary Meeting, not of Doctrine but of inward Experience, in which both Men found themſelves confronted by the ſame dreadful Queſtion: namely, whether the inward Structures which hold the World in its familiar Order might themſelves be altered, ſtrained, or undone.

Whether their converging Paths then diverged by deliberate Choice or by neceſſitous Fate, the Reader muſt judge. It is enough for the Editor to obſerve that, after this Period, Kant’s labours took the Form of a moſt rigorous and protective Critick, aimed at fortifying the Mind againſt the very Diſorders hinted in his Notes; while Swedenborg, weakened by Age and by inward Strife, appears to have turned his Viſions into a ſofter and more allegorical Strain, as though unwilling to tread again upon the trembling Ground he once approached too nearly.

\newpage
\part*{Part~III[0.5em]
On the Nature of the Boundary and the Neceſſity of the Veil}

\addcontentsline{toc}{part}{Part~III: On the Nature of the Boundary}

%%%%%%%%%%%%%%%%%%%%%%%%%%%%%%%%%%%%%%%%%%%%%%%%%%%%%%%%%%%%
%%%%%%%%%%%%%%%%%%%%%%%% CHAPTER IX %%%%%%%%%%%%%%%%%%%%%%%%%
%%%%%%%%%%%%%%%%%%%%%%%%%%%%%%%%%%%%%%%%%%%%%%%%%%%%%%%%%%%%

\chapter*{Chapter~IX[0.5em]
Wherein it is Argued that the Denial of Knowledge may be a Form of Mercy}

\addcontentsline{toc}{chapter}{Chapter~IX: The Denial of Knowledge as Mercy}

It hath often been ſaid that the firſt Office of Philoſophy is to enlarge the Bounds of human Knowledge; yet the Papers here collected ſuggeſt a moſt contrary Doctrine, namely, that the firſt and moſt neceſſary Office of a true Critick may be to confine thoſe Bounds, and to erect upon the Limits of Reaſon a Barrier more ſtrict than any erected by Ignorance alone.

For Ignorance breedeth mere Vacancy, and may be filled by the mild and natural Courſe of Experience; but a Knowledge that presseth beyond its rightful Province awakeneth Dangers not eaſily laid to Reſt. There is in theſe Papers a continual Theme of Precipice and Abyſs, as if the Mind, in its unwatched Wandering, might ſuddenly eſpy a Chasm beneath the ſeeming Firmneſs of the World.

In one paſſage of Kant’s private Notes, the following Confession appeareth:

\begin{quote}
``I have perceived the trembling of the Frame upon which all Appearance dependeth; and the Sight was more dreadful than any Material Horror. For the Ground of things withdrew like a Curtain, and I ſtood on Nothing.''
\end{quote}

And in another, evidently written under a moſt inward Perturbation:

\begin{quote}
``Were the Mind to behold the Thing in itſelf without Preparation, it would be rent, as a Picture torn from its Frame. Better a Boundary well-drawn than a View too vaſt for the Eye to endure.''
\end{quote}

Such Words, coming from a Man of that Temper and Diſcipline, do not convey the idle Terror of a Fancy ungoverned, but ſeem rather the Teſtimony of one who hath looked upon a Truth he wiſhed­—for the Sake of the Mind’s own Stability—not to behold again.

Swedenborg, from his opposite Vantage, reacheth a like Concluſion, though in Language more ſuited to his visionary Genius. He declareth that the Viſions granted him `were tempered according to the Weakneſs of the Mortal Frame,'' and that `to remove the Veil all at once would be to unmake the Creature who beholdeth it.''

Thus we perceive a singular Convergence of the two, wherein the Critick and the Seer alike affirm that the Limit of Knowledge is not a Priſon but a Mercy; not an Impediment, but a Shield drawn over the trembling Lattice that ſupporteth our Experience.

%%%%%%%%%%%%%%%%%%%%%%%%%%%%%%%%%%%%%%%%%%%%%%%%%%%%%%%%%%%%
%%%%%%%%%%%%%%%%%%%%%%%% CHAPTER X %%%%%%%%%%%%%%%%%%%%%%%%%%
%%%%%%%%%%%%%%%%%%%%%%%%%%%%%%%%%%%%%%%%%%%%%%%%%%%%%%%%%%%%

\chapter*{Chapter~X[0.5em]
On the Veil, and of its Neceſſity for the Stability of Experience}

\addcontentsline{toc}{chapter}{Chapter~X: On the Veil}

It is not uncommon among the Writers of Antiquity to ſpeak of a Veil that hideth the Godhead from mortal Sight; but in the Documents here examined, the Veil appeareth in a more intrinſick and critical Light. It is not merely a Cover HIDING some Higher Being, but rather an inward Law of the Mind itſelf, without which no orderly Experience would be poſſible.

In a detached Leaf found among Kant’s Effects, one readeth:

\begin{quote}
``The Veil is the Form of all Appearance; it is the Weaving of Space and Time themſelves. To pierce it is not to come nearer to Truth, but to unweave that by which Truth is apprehended.''
\end{quote}

Another Note, written in a more hurried Hand, declareth:

\begin{quote}
``The Veil hath two Sides: the one turned toward the World, the other toward the Abyss. To look beyond it, even for a Moment, is to hazard the loſs of both.''
\end{quote}

Swedenborg’s Journals, though differing widely in Doctrine, corroborate this fundamental Idea. More than once he reporteth that the Viſions vouchſafed unto him were `mitigated by Divine Mercy,'' for that `the naked Order of the World, as it ſtandeth before the Angelick Underſtanding, would overwhelm the mortal Senſe.''

From the Concurrence of theſe two Teſtimonies, the Editor inferreth that the Veil is not a Thing placed arbitrarily before the Mind, but a neceſſary Law of its Operation. It is the Frame, or the Boundary, by which the Chaos of manifold Senſations is given Shape and Order; and if that Frame ſhould ſhift, fray, or momentarily fail, the Mind would find itſelf confronted with a torrent of Unconditioned Being, againſt which no Creature may hope to ſtand.

Thus the Veil is at once our Limitation and our Refuge.

%%%%%%%%%%%%%%%%%%%%%%%%%%%%%%%%%%%%%%%%%%%%%%%%%%%%%%%%%%%%
%%%%%%%%%%%%%%%%%%%%%%%% CHAPTER XI %%%%%%%%%%%%%%%%%%%%%%%%%
%%%%%%%%%%%%%%%%%%%%%%%%%%%%%%%%%%%%%%%%%%%%%%%%%%%%%%%%%%%%

\chapter*{Chapter~XI[0.5em]
On the World as Stabilized Appearance}

\addcontentsline{toc}{chapter}{Chapter~XI: On Stabilized Appearance}

In contemplating the various Notes and ſecret Memoranda of theſe Years, one is ſtruck by a recurring Conceit: the World is not given all at once, nor in its whole Extent, but is continually maintained in its Order by the inward Activity of the Mind, operating under Laws which are themſelves unknown in their Source.

It hath been the Cuſtom of many Philoſophers to imagine the World as a vaſt Machine, whoſe Operations proceed independently of the Spectator; yet the Teſtimonies of theſe Papers ſuggeſt a more delicate Arrangement. In one fragment, Kant writeth:

\begin{quote}
``The World is appearance holding firm. Its Firmneſs is a Labour not of things but of the Mind's own Form. Should that Labour falter, even a Stone would forget how to fall.’’
\end{quote}

Another Margin containeth the extraordinary Note:

\begin{quote}
``Objects do not ſtand; they are held. Colours do not reſide; they are ſuſpended. The World is a Quiet maintained againſt a ceaseless Murmur beneath.''
\end{quote}

Theſe Expreſſions, though poetical in Sound, accord not ill with the Critick’s later Doctrine, in which Reaſon is repreſented as the Arbiter of the Poſſibility of Experience, and the Categories as the Conditions under which alone Objects may be given.

Swedenborg, though uſing a wholly different Scheme of Thought, ſpeaketh in his Journals of `the continual Tendency of things to fall back into the Abyss from which they were drawn,'' and of `the ceaseless ſuperintendence by which Nature is kept from lapsing into its firſt Unſhape.''

Thus both Men, approaching the Matter from oppoſite Extremes, conſpire in the Notion that the World as we apprehend it is a Kind of Stabilized Image or Conſtruction, maintained by Laws more inward than the Objects which they govern. And herein lieth the dreadful Poſſibility that theſe Laws may be ſubject to Strain, Diſturbance, or Interruption.

%%%%%%%%%%%%%%%%%%%%%%%%%%%%%%%%%%%%%%%%%%%%%%%%%%%%%%%%%%%%
%%%%%%%%%%%%%%%%%%%%%%%% CHAPTER XII %%%%%%%%%%%%%%%%%%%%%%%%
%%%%%%%%%%%%%%%%%%%%%%%%%%%%%%%%%%%%%%%%%%%%%%%%%%%%%%%%%%%%

\chapter*{Chapter~XII[0.5em]
Concerning the Dangerous Proximity of Explanation}

\addcontentsline{toc}{chapter}{Chapter~XII: On the Danger of Explanation}

If the preceding Chapters have in any manner prepared the Reader for that which followeth, it will not be thought extravagant when the Editor now ventureth to propoſe that in the Affairs of the Mind, there is a Point at which Explanation itſelf becometh dangerous.

For it is not true, as the vulgar Suppoſe, that all Diſcovery is beneficent, or that the Light of Knowledge may be extended without Limit. There are Truths which, though certain in themſelves, are too vaſt or too inward for the Mortal Frame to bear.

Among the Papers of Kant, one findeth this moſt remarkable and fearful Line:

\begin{quote}
``Some Explanations do not explain; they undo.''
\end{quote}

And in another Note:

\begin{quote}
``To know too nearly why things appear is to endanger the Power by which they appear.''
\end{quote}

In the Journals of Swedenborg we encounter a Parable of this Nature:

\begin{quote}
``There was a Man who demanded to see the Thread by which his Soul was joined to the World. And the Angels ſaid: The Thread is finer than the Breath of a Candle. Should we ſhew it, it would break. And the Man, not believing, begged again; and they ſhewed it. And in the Moment he beheld it, the Thread was cut, and he fell as one dead.''
\end{quote}

Theſe Teſtimonies, taken together, admoniſh us that the Boundary which separateth Appearance from its hidden Ground is not to be lightly transgreſſed. To explain too nearly may not merely diſturb the Mind that attempteth it, but may alſo diſcompoſe the Order of Experience itſelf.

It is perhaps for this Reaſon that Kant, after theſe Dark Years, ſet himſelf with ſuch iron Reſolution to the Task of erecting a Critick that would foreſtall any future Mind from falling into the ſame inward Diſorder. His Labours were not, as ſome have jested, the Idle Scruples of a Cloiſtered Scholar, but the Defenſive Works of a Man who had ſeen too much, and who would fain erect a Fortress whereby the Mind might be guarded from its own Ambition.

\newpage

\part*{Part~IV[0.5em]
Manuſcripts of Immanuel~Kant (Nachlaſs~K–42)}

\addcontentsline{toc}{part}{Part~IV: Manuſcripts of Kant}

%%%%%%%%%%%%%%%%%%%%%%%%%%%%%%%%%%%%%%%%%%%%%%%%%%%%%%%%%%%%
%%%%%%%%%%%%%%%%%%%%%%%% CHAPTER XIII %%%%%%%%%%%%%%%%%%%%%%%
%%%%%%%%%%%%%%%%%%%%%%%%%%%%%%%%%%%%%%%%%%%%%%%%%%%%%%%%%%%%

\chapter*{Chapter~XIII[0.5em]
Notes Concerning the Event of the Year~1769}

\addcontentsline{toc}{chapter}{Chapter~XIII: Notes Concerning the Event of 1769}

It is a Matter of ſingular Curioſity that among the private Leaves of Kant, preſerved in a Packet marked only with the Numeral~XLII, there ſhould be found ſo many cryptic Paſſages relating to a ſingle Winter. Though the Editor maketh no Claim to understand their full Import, the inward Tone of the Writings betrayeth a Mind confronting a moſt profound Diſturbance.

One Leaf, written in a fine but ſomewhat trembling Hand, preſenteth the following:

\begin{quote}
``The Year hath turned, and with it my Perceptions.
The Shadow of the Clock-tower ſeemed, this Evening, to lag a Hand-breadth behind the Tower itſelf. A Moment only; yet in that Moment the Order of the World faltered. I would have attributed it to the Eye’s wearineſs, had not the freezing of the Air ſeemed to wait upon the Shadow, as though Time were caught in a Knot.''
\end{quote}

Another Fragment, clearly from the ſame Night:

\begin{quote}
``If this be a Deluſion, it is of ſuch a Kind as I have not known. The Lamps along the Kalte Straße burned with a uniform Flame, yet my Steps made no Echo. The Silence preſſed againſt me like a Hand. I felt as though I were walking through Appearance newly made, not yet ſettled.''
\end{quote}

And a third:

\begin{quote}
``The Diagram recurred to me as I ſlept, tracing itſelf upon the inward Eye. Four Lines, meeting at Angles not euclidean, yet not without Law. If this be a Fancy, why doth my Hand attempt to draw it even as I wake?''
\end{quote}

It is apparent that theſe Notes were never intended for publick View, nor even for the more intimate Circle of Kant’s Students. Their Tone is urgent, confiding, at Moments confessional, as if the Writer feared he had perceived ſomething he ought not to have ſeen.

%%%%%%%%%%%%%%%%%%%%%%%%%%%%%%%%%%%%%%%%%%%%%%%%%%%%%%%%%%%%
%%%%%%%%%%%%%%%%%%%%%%%% CHAPTER XIV %%%%%%%%%%%%%%%%%%%%%%%%
%%%%%%%%%%%%%%%%%%%%%%%%%%%%%%%%%%%%%%%%%%%%%%%%%%%%%%%%%%%%

\chapter*{Chapter~XIV[0.5em]
Fragments on Shadows and their Obedience}

\addcontentsline{toc}{chapter}{Chapter~XIV: Shadows and their Obedience}

No Fragments in the whole Packet XLII are more diſquieting than thoſe which concern Shadows. For in theſe Paſſages, normally brief, abrupt, and bereft of all the graces of Style, we encounter the ſudden Introduction of a Notion which Kant never afterwards employed in his publick Writings, yet which appears to have gripped his Mind with peculiar Force.

One Fragment runneth thus:

\begin{quote}
``The Shadows obeyed Me this Morning. I walked before the Sun; yet the Shadow hesitated, then followed. A Delay of perhaps a twelfth-part of a Second. The Eye might err, but the Body felt it.''
\end{quote}

Another Leaf recordeth:

\begin{quote}
``When I extended my Hand toward the Stove, the Shadow did not extend. It ſtood as it had ſtood. I drew back in Terror. The Flame flickered then as if in acknowledgement. My Shadow then corrected itſelf, with a Motion too eaſy to be believed.''
\end{quote}

And in a Line written in a cramped Hand along the margin:

\begin{quote}
``If the Shadow be the Sign of Succession, what then if it rebel? If the ſign that proclaims Time ſhould diverge from the Body whose Time it marketh, hath not the Order of the World begun to ſhift?''
\end{quote}

It is notable that this ſtrain of Thought convergeſth—though in a wholly different Idiom—with a Paſſage found in the later Journal of Swedenborg, wherein he affirmeth that ``the inverted Forms move before their Objects, ſignifying that Cauſality hath loſt its Direction.''

The Editor, though averſe to all idle Speculation, cannot but remark that theſe Paſſages are mutually illuminative, and point toward a claſſ of inward Experience of which no precedent is found in the Writings of the Enlightened.

%%%%%%%%%%%%%%%%%%%%%%%%%%%%%%%%%%%%%%%%%%%%%%%%%%%%%%%%%%%%
%%%%%%%%%%%%%%%%%%%%%%%% CHAPTER XV %%%%%%%%%%%%%%%%%%%%%%%%%
%%%%%%%%%%%%%%%%%%%%%%%%%%%%%%%%%%%%%%%%%%%%%%%%%%%%%%%%%%%%

\chapter*{Chapter~XV[0.5em]
Obſervations upon the Inverſion of the City beneath the River}

\addcontentsline{toc}{chapter}{Chapter~XV: On the Inverſion of the City}

Among the moſt extraordinary Leaves of Packet XLII is one which appeareth to deſcribe a Viſion or inward Perception experienced upon the frozen Pregel during the Winter of that Year. It is in a Hand more hurried than Kant's usual Manuſcript, and certain Words are ſtricken through or repeated, as if the Writer faltered between Fear and Precision.

The Fragment runneth thus:

\begin{quote}
``I walked at Evening upon the Pregel, the Ice bearing firm. The Cold was sharp, but the Air ſtill. I felt drawn to the Water's Middle, though I knew not why.

There, before Me, the Ice thinned, and a Circle of Water opened, perfect as if traced by Compass. Beneath the Surface ſhone a City—not ours, yet ours—the Streets inverted, the Spires pointing into the Depths. A Bell tolled without Sound; the Air in My Breaſt trembled as at an inward Thunder.

The City turned then, as if beholding Me. I fell to My Knees, yet I could not break My Gaze.

When the Water cloſed again, My Reflection remained; but it did not kneel. It studied Me. The Mouth moved, yet Mine did not.

I fled, yet My Shadow did not flee at once.''
\end{quote}

Whatſoever the Cauſe, this Paſſage revealeth an inward Diſorder of a moſt rare and terrifying Nature. It is followed by another Fragment, evidently written upon returning to his Study:

\begin{quote}
``If Perception itſelf be invertible, then the World I apprehend is but one Branch of a larger Tree. I dare not purſue this Thought; it will undo my Reaſon.’’
\end{quote}

Thus the Inverſion of the City beneath the River, whether viſionary or perceptual, muſt not be regarded as a mere Fancy but as a ſignificant Moment in the Formation of Kant’s later Conviction concerning the Limits of Knowledge.

%%%%%%%%%%%%%%%%%%%%%%%%%%%%%%%%%%%%%%%%%%%%%%%%%%%%%%%%%%%%
%%%%%%%%%%%%%%%%%%%%%%%% CHAPTER XVI %%%%%%%%%%%%%%%%%%%%%%%%
%%%%%%%%%%%%%%%%%%%%%%%%%%%%%%%%%%%%%%%%%%%%%%%%%%%%%%%%%%%%

\chapter*{Chapter~XVI[0.5em]
The Command: “Write the Boundary”}

\addcontentsline{toc}{chapter}{Chapter~XVI: The Command to Write the Boundary}

Of all the Leaves in Packet XLII, none hath excited more Speculation among later Commentators than the Page upon which Kant recounteth a moſt alarming Experience, marked by the Appearance of a Line of Writing that he affirmeth was not of his own Hand.

The Paſſage readeth thus:

\begin{quote}
``Tonight the Board wrote again. The Chalk lay untouched upon the Tray; yet a Line appeared, carved into the Slate with a silent Force.

WRITE THE BOUNDARY.

I stood frozen. The Letters were deep; the Hand that made them was unknown.

I cried out: `I am trying.'

The Chalk then traced the Diagram, the four Lines bending as if the Board itſelf possessed Will.

When I placed my Hand upon the Slate, the Stone throbbed with a Pulse not mine.

I know not what the Command intendeth, yet its Insistence groweth: Write the Boundary, else the World will not hold.''
\end{quote}

Another Note, found folded into the ſame Page, addeth:

\begin{quote}
``This is not a Dictate of Reaſon, nor of any Faculty known to me. It is the Voice of the Order that upholdeth Experience. If I do not obey, the Frame may tremble. If I obey, I may never again think freely.

Yet I must choose.’’
\end{quote}

The Editor hazardeth no definite Interpretation of theſe Paſſages, but obſerveth only that the Command, whatſoever its Origin, appeareth later in Kant’s mature Writings transformed into the Doctrine of the Limits of Reaſon, and the neceſſary Diſtinction between Phenomenon and Noumenon.

%%%%%%%%%%%%%%%%%%%%%%%%%%%%%%%%%%%%%%%%%%%%%%%%%%%%%%%%%%%%
%%%%%%%%%%%%%%%%%%%%%%%% CHAPTER XVII %%%%%%%%%%%%%%%%%%%%%%%
%%%%%%%%%%%%%%%%%%%%%%%%%%%%%%%%%%%%%%%%%%%%%%%%%%%%%%%%%%%%

\chapter*{Chapter~XVII[0.5em]
Two Manuſcripts: The Mask and the Truth}

\addcontentsline{toc}{chapter}{Chapter~XVII: Two Manuſcripts: The Mask and the Truth}

In sorting the Contents of Packet XLII, the Editor found evidence—nowhere acknowledged in Kant's publick Writings—that during the Years 1767–1768 he compoſed not one but two extended Treatiſes. One of theſe, polished and tempered for the publick Eye, later appeared under the Title \textit{Träume eines Geistersehers}. The other, rough, unbound, and filled with Corrections and marginal Exclama­tions, containeth the more fearful Conceptions that were afterwards tranſmuted into the \textit{Critique of Pure Reason}.

On a loose Sheet affixed to the latter, Kant writeth:

\begin{quote}
``I muſt hide the Truth beneath a Mask. The World is not prepared. Let the publick laugh at the Spirit-Seer, and take comfort in believing that I have allied myſelf with their Diſdain.

But the greater Work, which containeth the Boundary, muſt be forged in Silence.’’
\end{quote}

Another Leaf, evidently written after a Night of inward Conflict:

\begin{quote}
``The Mask will keep them ſafe. The Truth will keep the World from breaking. I muſt publiſh the one, and ſeal the other.’’
\end{quote}

Thus the two Manuſcripts—one to appeaſe, one to preſerve—became the dual Fruits of that extraordinary Period, and their Divergence is among the moſt remarkable Facts in the Hiſtory of the Human Mind.

%%%%%%%%%%%%%%%%%%%%%%%%%%%%%%%%%%%%%%%%%%%%%%%%%%%%%%%%%%%%
%%%%%%%%%%%%%%%%%%%%%%%% CHAPTER XVIII %%%%%%%%%%%%%%%%%%%%%%
%%%%%%%%%%%%%%%%%%%%%%%%%%%%%%%%%%%%%%%%%%%%%%%%%%%%%%%%%%%%

\chapter*{Chapter~XVIII[0.5em]
Late Reflections on Time and the Price Paid}

\addcontentsline{toc}{chapter}{Chapter~XVIII: Late Reflections and the Price Paid}

The final Leaves of Packet XLII are of a Tone markedly quieter than thoſe preceding. The Fury of Perception hath abated; yet its Memory lingereth, and the Writings bear the Stamp of a Man who hath ſeen too deeply into the inward Frame of Experience, and who beareth the Weight of that Knowledge with a grave Reſolution.

One such Leaf recordeth:

\begin{quote}
``This Night the Shadows obeyed. This Night the Board was ſilent. The Diagram did not appear. Perhaps the Boundary is firm now.

Yet the Price is heavy: I ſhall never again truſt the Simplicity of Experience. Everywhere I perceive the Labour by which the Mind upholdeth the World.’’
\end{quote}

Another, evidently from his later Years:

\begin{quote}
``The Frame held. The World kept its Quiet. But I perceive now that this Quiet is not Natural, but Won. And what is Won may be loſt, if the Mind forgetteth its Limits.

Thus I muſt teach the Limits, though I would gladly have taught the Freedom beyond them.’’
\end{quote}

The final Fragment of Packet XLII readeth:

\begin{quote}
``Let them call me narrow, ſevere, cautious. They know not the Abyss I have ſeen. If the Boundary holdeth, the World may continue in its gentle Appearance.

This is enough.’’
\end{quote}

\newpage

\part*{Part~V\[0.5em]
The Laſt Journal of Emanuel~Swedenborg (MDCCLXX–MDCCLXXII)}

\addcontentsline{toc}{part}{Part~V: Laſt Journal of Swedenborg}

%%%%%%%%%%%%%%%%%%%%%%%%%%%%%%%%%%%%%%%%%%%%%%%%%%%%%%%%%%%%
%%%%%%%%%%%%%%%%%%%%%%%% CHAPTER XIX %%%%%%%%%%%%%%%%%%%%%%%%
%%%%%%%%%%%%%%%%%%%%%%%%%%%%%%%%%%%%%%%%%%%%%%%%%%%%%%%%%%%%

\chapter*{Chapter~XIX\[0.5em]
Early Entries on Space, Light, and Meaning}

\addcontentsline{toc}{chapter}{Chapter~XIX: Early Entries}

The Journal of the venerable Swedenborg, covering the Years~1770 to 1772, was found in a small wooden Casket bound with Iron. The Leaves are worn, many Pages blotted by Moisture, yet the Hand is clear, and the inward Spirit of the Writer is discernible in every Line.

The earlier Entries of this Period reveal an undercurrent of Anxiety, as though the aged Seer perceived the gradual dimming of his Viſions, or the withdrawal of that inward Light which he had believed Divine.

One Entry, dated the Third Day of January~1770, readeth:

\begin{quote}
``I walked in the Garden behind My Lodging, though the Snow lay deep. The Air was ſtill, the Light thin and of a cold Purity. As I ſurveyed the Morning, I perceived faintly the Threads that bind Space into its Juſt Proportion. They were dim, as though veil'd in a new Manner.

Where once the Lines were bright, they now flicker. And I am left to wonder whether my Sight dimmeth, or whether the World bluſheth with ſome inward Shame.’’
\end{quote}

A later Entry from the ſame Month declareth:

\begin{quote}
``The Light is changed. It falleth upon Objects without that inward Meaning which once pierced Me. I fear that the Kingdoms I beheld in Viſion withdraw their Radiance, as though preparing Me for Silence.

Yet, the World trembleth. I have ſeen Windows reflect the City not as it is, but as it might be if the inward Frame were to bend.’’
\end{quote}

Swedenborg’s habitual Strain toward Allegory in theſe Lines ſhould not obscure the poignant Sense that he confronted a waning of his viſionary Faculty, coupled with a fear—uncharacteristic for his buoyant Temper—that the inward Structures he believed himſelf to perceive were themſelves under a kind of Diſturbance.

%%%%%%%%%%%%%%%%%%%%%%%%%%%%%%%%%%%%%%%%%%%%%%%%%%%%%%%%%%%%
%%%%%%%%%%%%%%%%%%%%%%%% CHAPTER XX %%%%%%%%%%%%%%%%%%%%%%%%%
%%%%%%%%%%%%%%%%%%%%%%%%%%%%%%%%%%%%%%%%%%%%%%%%%%%%%%%%%%%%

\chapter*{Chapter~XX\[0.5em]
On the Inverſed City in Dream and in Waking}

\addcontentsline{toc}{chapter}{Chapter~XX: Inverſed City}

In the midſt of the Journal’s earlier Entries, we encounter a ſtriking Series of Paſſages concerning a Viſion (or Dream) of an Inverſed City. Though Swedenborg had, throughout his Life, described many Viſions of the Spiritual Kingdoms, this particular Conception is of an altogether different Character.

From an Entry dated the Tenth of March~1770:

\begin{quote}
``In the Night I dreamt again of the Inverſed City, which I once glimpsed ſome Years ago beneath a frozen Lake in my own Land. But now the Viſion was clearer.

Streets falling downward into Depths. Gardens ſpread beneath Me, their Trees pointing toward the Earth’s Centre. Bells that tolled not by Vibration, but by the Shiver of the inward Frame.

I felt the City behold Me. And I knew then that its Inversion was not of Place, but of Cauſe; that its Order ran backward from the World we know.’’
\end{quote}

The Entry following on the next Leaf, written at Dawn, continueth:

\begin{quote}
``I walked in Waking toward the Canal. There, in a Sheen of Water not yet frozen, I perceived the ſame Image faintly; not so bright as in the Dream, but undeniably present.

It ſeemed as though the World’s outward Face faltered, and another Arrangement ſtood beneath it, waiting to be acknowledged.

Is this the Sign that the inward Lattice quivers? If so, then we who perceive it are in peril, for we have ſeen the Frame without the Tapestry.’’
\end{quote}

This Conception of an Inverſed City—both dreamt and perceived—mirrors with remarkable exactitude the Fragment recorded in Kant’s own Notes. The Editor remarketh, without Comment, that both Men appear to have encountered the ſame inward Structure, though they interpreted it through ſingularly different Faculties.

%%%%%%%%%%%%%%%%%%%%%%%%%%%%%%%%%%%%%%%%%%%%%%%%%%%%%%%%%%%%
%%%%%%%%%%%%%%%%%%%%%%%% CHAPTER XXI %%%%%%%%%%%%%%%%%%%%%%%%
%%%%%%%%%%%%%%%%%%%%%%%%%%%%%%%%%%%%%%%%%%%%%%%%%%%%%%%%%%%%

\chapter*{Chapter~XXI\[0.5em]
On Kant: His Strength and His Terror}

\addcontentsline{toc}{chapter}{Chapter~XXI: On Kant}

In the Journal’s Mid-Year Entries, we find a curious Shift in Tone, wherein Swedenborg addreſſeth his Thoughts to a ``certain Man of Königsberg,'' whom the Editor taketh to be none other than Professor~Kant.

An Entry dated the Twenty-ſecond of July~1771 declareth:

\begin{quote}
``The Man of Königsberg hath approached the Boundary which I have walked theſe many Years. Yet he cometh not as a Seer, but as an Architect of Reaſon.

He trembleth. His Shadow delayeth. His Perceptions quiver as one who draweth too near a Furnace.

But he hath a Strength I do not. He could draw the Line that neither I nor any Angel dare draught. He could build the Wall that concealeth what would undo the World.’’
\end{quote}

A later Entry, written upon his return from a Visit to Stockholm, readeth:

\begin{quote}
``He feareth; yet he is not overwhelmed. He would deny even his own Eyes if thereby he might protect the Order by which Men live.

His Fate is heavier than mine. For I have ſeen the Kingdoms beyond the Veil and have rejoiced. But he hath ſeen the trembling of the Veil itſelf, and hath felt the World’s Dependence upon its Strength.

I honour him. I pity him.’’
\end{quote}

Theſe Entries reveal a moſt poignant Admiration, mingled with a Recognition that Kant’s inward Burden was of a different Order from Swedenborg’s; for where the one sought to tranſcend the Veil, the other sought to preſerve it.

%%%%%%%%%%%%%%%%%%%%%%%%%%%%%%%%%%%%%%%%%%%%%%%%%%%%%%%%%%%%
%%%%%%%%%%%%%%%%%%%%%%%% CHAPTER XXII %%%%%%%%%%%%%%%%%%%%%%%
%%%%%%%%%%%%%%%%%%%%%%%%%%%%%%%%%%%%%%%%%%%%%%%%%%%%%%%%%%%%

\chapter*{Chapter~XXII\[0.5em]
Approaching Death and the Thinning Veil}

\addcontentsline{toc}{chapter}{Chapter~XXII: Approaching Death}

As the Journal proceedeth into the late Months of 1771 and the early Spring of 1772, a marvellous Tranſition taketh place in the Style of the Entries. The Language groweth sparse, the Imagery more inward, the Tone more sombre. Swedenborg’s viſionary Power, which in earlier Years was bright and exuberant, is here tempered by a ſolemn Recognition that his Strength waneth, and that the Worlds he beheld ſtand now at a greater Remove.

An Entry of November~1771:

\begin{quote}
``The Veil groweth firmer for me. Not by Strength, but by Weakneſs. The Eyes dim; the inward Light fainteth. The Kingdoms which once ſhone around me now appear but as faint Embers.

Yet the World trembleth ſtill; I ſee it in the Tremor of Shadows and in the brief Faltering of Cauſe.''
\end{quote}

A Paſſage from the Eleventh Day of February~1772 is particularly moving:

\begin{quote}
``I feel the Approach of that Departure which no Man may eſcape. The inward Lattice groweth simpler; the Threads are few. I ſee no more the manifold Orders of Spirits, but only a certain Quiver at the Root of Appearance.

Perhaps Death is but the Soul’s Detachment from the Lattice; perhaps Life is but a fitting into its Threads.

I feel the Threads loosening.’’
\end{quote}

These Paſſages, though hallowed by the melancholy of Age, reveal a remarkable Serenity in the Face of his weakening Faculty. The trembling of the inward Lattice, which had once filled him with Awe, now ſeemeth but a Faint Murmur at the Edge of Departure.

%%%%%%%%%%%%%%%%%%%%%%%%%%%%%%%%%%%%%%%%%%%%%%%%%%%%%%%%%%%%
%%%%%%%%%%%%%%%%%%%%%%%% CHAPTER XXIII %%%%%%%%%%%%%%%%%%%%%%
%%%%%%%%%%%%%%%%%%%%%%%%%%%%%%%%%%%%%%%%%%%%%%%%%%%%%%%%%%%%

\chapter*{Chapter~XXIII\[0.5em]
The Final Entry: ``The Witneſs may Reſt’’}

\addcontentsline{toc}{chapter}{Chapter~XXIII: Final Entry}

The laſt Entry in Swedenborg’s Journal, dated the Twenty-ninth Day of March~1772, and believed by many to have been penned within Hours of his Departure, is brief but of a moſt profound Simplicity.

\begin{quote}
``This Day I was called Home.

I have ſeen the Lattice tremble; I have ſeen it ſtand. The Man of Königsberg hath done what I could not: he hath drawn the Boundary that hideth the inward Frame.

Tell him: the Lattice held. The Price is paid.

The Witneſs may now reſt.’’
\end{quote}

Thus endeth the Laſt Journal of Emanuel Swedenborg.

\newpage
\part*{Part~VI[0.5em]
A Scholarly Appendix Concerning Parallel Paſſages, Chronology, and Diverſe Interpretations}

\addcontentsline{toc}{part}{Part~VI: Scholarly Appendix}

%%%%%%%%%%%%%%%%%%%%%%%%%%%%%%%%%%%%%%%%%%%%%%%%%%%%%%%%%%%%
%%%%%%%%%%%%%%%%%%%%%%%% CHAPTER XXIV %%%%%%%%%%%%%%%%%%%%%%%
%%%%%%%%%%%%%%%%%%%%%%%%%%%%%%%%%%%%%%%%%%%%%%%%%%%%%%%%%%%%

\chapter*{Chapter~XXIV[0.5em]
Parallel Paſſages in Kant and Swedenborg, with Remarks}

\addcontentsline{toc}{chapter}{Chapter~XXIV: Parallel Paſſages}

The Editor, having carefully collated the Manuſcripts of Kant (Nachlaſs~K–42) with the Laſt Journal of Swedenborg (MDCCLXX–MDCCLXXII), judged it profitable to preſent certain notable Paſſages in parallel, with no aim other than to aid the Reader’s Contemplation. No Concluſions are forced; the Texts are left to illumine one another by their natural Affinity.

\medskip

\noindent\textit{1. Concerning the trembling of the inward Frame}

\begin{quote}
\textbf{Kant (Nachlaſs XLII):}
``I perceived the trembling of the Frame upon which all Appearance dependeth, and the Sight was dreadful.''

\textbf{Swedenborg (Journal, Jan.~1770):}
``The inward Threads that bind Space into Order flicker’d, as though the World bluſh’d with ſome inward Shame.’’
\end{quote}

Both Paſſages, though couched in different Idioms, expreſs a conception of an inward Architecture upon which all Appearance dependeth, and which may tremble or falter.

\medskip

\noindent\textit{2. On Inverſion of Place}

\begin{quote}
\textbf{Kant (Nachlaſs):}
``The City beneath the River turned as if beholding Me; its Order ran contrary to the World I know.’’

\textbf{Swedenborg (Journal, March~1770):}
``I ſaw the Inverſed City again, its Streets falling into Depths, its Cauſe running backward from the World.’’
\end{quote}

The Agreement here is moſt ſtriking, given their independent Provenance.

\medskip

\noindent\textit{3. On Shadows and Succession}

\begin{quote}
\textbf{Kant:}
``The Shadow delayed; then corrected itſelf, as if in acknowledgement.’’

\textbf{Swedenborg:}
``I ſaw Forms move before their Objects, ſignifying Cauſality unmoored.’’
\end{quote}

Both Teſtimonies touch upon a Diſturbance of the Relation between Object and its Shadow, a Sign of temporal Diſorder.

\medskip

\noindent\textit{4. On the Neceſſity of the Boundary}

\begin{quote}
\textbf{Kant:}
``Write the Boundary, elſe the World will not hold.’’

\textbf{Swedenborg:}
``He alone could draw the Line that concealeth what would undo the World.’’
\end{quote}

Though expreſſed differently, both Paſſages suggest that the Stability of Experience dependeth upon a certain inward Limit or Boundary.

\medskip

\noindent Such Convergences, taken together, indicate that the inward Experience of both Men during theſe Years bore upon the ſame Metaphyſical Problem, though their Tempero and Vocations led them to very different Expressions.

%%%%%%%%%%%%%%%%%%%%%%%%%%%%%%%%%%%%%%%%%%%%%%%%%%%%%%%%%%%%
%%%%%%%%%%%%%%%%%%%%%%%% CHAPTER XXV %%%%%%%%%%%%%%%%%%%%%%%%
%%%%%%%%%%%%%%%%%%%%%%%%%%%%%%%%%%%%%%%%%%%%%%%%%%%%%%%%%%%%

\chapter*{Chapter~XXV[0.5em]
A Chronology of the Anomalous Year}

\addcontentsline{toc}{chapter}{Chapter~XXV: Chronology}

The following Chronology, though imperfect in its Reliance upon fragmentary Teſtimonies, preſents, as nearly as the Evidence permitteth, the Principal Events of the ſingular Year~1769 and its immediate Aftermath. Theſe Dates are offered with academic Caution; the Reader muſt not deem them certain where the Nature of the Sources is oblique.

\medskip

\noindent\textbf{Winter of 1768–1769.}
Several early Notes in Kant’s Hand indicate Perceptual Diſturbances: trembling Shadows, the ſilent bending of Candle-flames, and the firſt Appearance of the Diagram which recurreth throughout Packet~XLII.

\medskip

\noindent\textbf{January 1769.}
Kant recordeth the firſt Experience of the World’s Momentary Silence: ``My Steps made no Echo; the Air preſs’d like a Hand.''

\medskip

\noindent\textbf{February 1769.}
Reports of Swedenborg’s extraordinary Predictions circulate in Königsberg; Kant takes moſt pointed Intereſt, though outwardly diſdainful.

\medskip

\noindent\textbf{March–April 1769.}
Kant’s Notes refer to ``the Inverſion beneath the Pregel,'' likely the Viſion of the Inverſed City.

\medskip

\noindent\textbf{June 1769.}
First Indication of the Command: ``WRITE THE BOUNDARY.''

\medskip

\noindent\textbf{Late 1769.}
Correspondence between the two Men (see Part~VII) ſuggeſts a brief, ſubtle Interchange of Ideas.

\medskip

\noindent\textbf{Early 1770.}
Swedenborg’s Journal reports a waning of his viſionary Faculty and renewed Anxiety over the trembling of the inward Lattice.

\medskip

\noindent\textbf{1771.}
Both Men refer indirectly to the other in their Notes, with mutual Admiration tinged with Pity.

\medskip

\noindent\textbf{March 1772.}
Swedenborg’s laſt Entry: ``Tell him: the Lattice held. The Price is paid.''

\medskip

\noindent This Chronology, though fragmentary, exceedeth by far the known Synchrony between the two Lives, and revealeth a Convergence of inward Experience not explicable by the Ordinary Courſe of Events.

%%%%%%%%%%%%%%%%%%%%%%%%%%%%%%%%%%%%%%%%%%%%%%%%%%%%%%%%%%%%
%%%%%%%%%%%%%%%%%%%%%%%% CHAPTER XXVI %%%%%%%%%%%%%%%%%%%%%%%
%%%%%%%%%%%%%%%%%%%%%%%%%%%%%%%%%%%%%%%%%%%%%%%%%%%%%%%%%%%%

\chapter*{Chapter~XXVI[0.5em]
Croſs-References to the \textit{Critique of Pure Reaſon}}

\addcontentsline{toc}{chapter}{Chapter~XXVI: Croſs-References}

The Editor, being mindful that the Reader may wiſh to compare the private Notes of the preceding Parts with Kant’s later printed Writings, here offereth a few careful Croſs-References. The Intention is not to unveil any conjectured Source of the Critick’s mature Doctrine, but merely to note Curioſities of Thought which may indicate how early inward Perturbations were afterwards transmuted into Principles of the Critickal Philoſophy.

\medskip

\noindent\textbf{1. On Space and Time}
Kant’s private Notes speak of `the Lattice'' and `the Ribbon of Time.'' In the \textit{Transcendental Æſthetic}, he proclaimeth Space and Time to be pure Forms of the Intuition.

\medskip

\noindent\textbf{2. On the Stability of Appearance}
Private: `Objects are held, not standing.''  
Printed: `The Synthesis of Apprehension dependeth upon the Unity of Consciousneſs.’’

\medskip

\noindent\textbf{3. On the Limits of Knowledge}
Private: `Some Explanations undo.''  
Printed: `I had to deny Knowledge in order to make room for Faith.''

\medskip

\noindent\textbf{4. On the Thing-in-itſelf}
Private: `The Veil turneth two Ways; to behold the other is to hazard all.’’  
Printed: `The Ding-an-sich remaineth unknowable to the human Underſtanding.’’

\medskip

\noindent\textbf{5. On Shadows and Succession}
Private: ``The Shadow delayed.''
Printed: Though no ſuch Notion appeareth directly, the entire \textit{Schematism} may be seen as an attempt to reconcile the inner Temporal Order with the outward Appearance of Objects.

\medskip

\noindent Theſe Croſs-References, though far from exhauſtive, indicate how private inward Experiences may have contributed to the mature Formulation of a Doctrine which, in its outward Severity, concealeth its inward Origin.

%%%%%%%%%%%%%%%%%%%%%%%%%%%%%%%%%%%%%%%%%%%%%%%%%%%%%%%%%%%%
%%%%%%%%%%%%%%%%%%%%%%%% CHAPTER XXVII %%%%%%%%%%%%%%%%%%%%%%
%%%%%%%%%%%%%%%%%%%%%%%%%%%%%%%%%%%%%%%%%%%%%%%%%%%%%%%%%%%%

\chapter*{Chapter~XXVII[0.5em]
Notes on Marginalia and Suppreſſed Leaves}

\addcontentsline{toc}{chapter}{Chapter~XXVII: Marginalia and Suppreſſed Leaves}

The Editor deemeth it neceſſary to record that several Leaves of Packet~XLII contain Marginalia that were evidently struck out by Kant himſelf. The Hand is his, but the Inks differ, suggesting that he returned to the Leaves in later Years and expunged Paſſages he judged too revealing or too incoherent.

One ſuch Note readeth:

\begin{quote}
``There is another Order beneath this Order. It presseth upward. It ſeeketh to break through the Veil. The Diagram is its fore-token.’’
\end{quote}

A thick Stroke of Ink attempts to obliterate the entire Line, yet the Words remain legible.

Another:

\begin{quote}
``If the World be but a Truce, then the Boundary is the Treaty. But who are the Contracting Parties?’’
\end{quote}

A third Margin containeth only the Words:

\begin{quote}
``The Order looketh back.’’
\end{quote}

These Leaves, though fragmentary and at times bordering upon the mystical, bear the Insight of a Mind striving to articulate an inward Apprehension of which no preceding Philoſopher had spoken.

The Editor leaves Judgment to the Reader.

%%%%%%%%%%%%%%%%%%%%%%%%%%%%%%%%%%%%%%%%%%%%%%%%%%%%%%%%%%%%
%%%%%%%%%%%%%%%%%%%%%%%% CHAPTER XXVIII %%%%%%%%%%%%%%%%%%%%%
%%%%%%%%%%%%%%%%%%%%%%%%%%%%%%%%%%%%%%%%%%%%%%%%%%%%%%%%%%%%

\chapter*{Chapter~XXVIII[0.5em]
Interpretive Difficulties of the Boundary-Conception}

\addcontentsline{toc}{chapter}{Chapter~XXVIII: Interpretive Difficulties}

No part of this Volume hath excited more Controverſy than the Conception of the ``Boundary.'' Some have argued that Kant employed the Term in a purely metaphorical Senſe, meaning only the Limitation of human Knowledge. Others contend that a more literal Meaning is here intended: a metaphyſical Boundary inherent in the inward Structure of Experience itſelf.

The Editor, while refraining from all Speculation beyond what the Texts warrant, remarketh that the Manuscripts present several genuine Difficulties:

\medskip

\noindent\textbf{1.} Kant’s repeated ſtatements that the Boundary must be \textit{written} (not merely recognized) imply that the Limit is not wholly given, but in some manner constructed or maintained.

\medskip

\noindent\textbf{2.} Several passaged affirme that without the Boundary, ``the World will not hold,'' which appear to refer not to the Limits of Inquiry but to the Stability of Appearance itſelf.

\medskip

\noindent\textbf{3.} Swedenborg’s parallel Assertions that the inward Lattice `quivereth,'' and that `the Price is paid,'' indicate that he perceived the Boundary as an active Structure with Consequences.

\medskip

\noindent\textbf{4.} Both Men, in their late Writings, refer to the Boundary not as a Theorem but as a Duty, a Task, or a Chore (``the Labour by which the World is upheld'').

\medskip

\noindent Thus the Interpretation of the Boundary cannot be reduced to a mere Intellectual Conception. It is presented here as a Metaphyſical Structure, an inward Frame, and as a Moral Duty incumbent upon the Mind that perceiveth too nearly the trembling of the World.

\newpage
\part*{Part~VII[0.5em]
The Secret Correſpondence (MDCCLXIX–MDCCLXXII)}

\addcontentsline{toc}{part}{Part~VII: Secret Correſpondence}

%%%%%%%%%%%%%%%%%%%%%%%%%%%%%%%%%%%%%%%%%%%%%%%%%%%%%%%%%%%%
%%%%%%%%%%%%%%%%%%%%%%%% CHAPTER XXIX %%%%%%%%%%%%%%%%%%%%%%%
%%%%%%%%%%%%%%%%%%%%%%%%%%%%%%%%%%%%%%%%%%%%%%%%%%%%%%%%%%%%

\chapter*{Chapter~XXIX[0.5em]
The Firſt Letter of Swedenborg, Delivered by Unknown Hand}

\addcontentsline{toc}{chapter}{Chapter~XXIX: The Firſt Letter}

The Document here tranſcribed is the Earliest known Piece of Correſpondence between the two extraordinary Minds who form the Subject of this Volume. The original Letter was found folded within Packet~XLII, bearing no Seal and no Signature, though the Hand is unmistakably that of Swedenborg. Its Date is uncertain, but internal Evidence suggesteth the late Winter of 1769.

\medskip

\noindent\textit{To the Learned Philoſopher dwelling in Königsberg,}

\begin{quote}
``It hath come to my inward Senſe that thou haſt perceived what ſhould not be perceived by any Man who hath not been prepared by long Diſcipline and inward Burning. The trembling of the outward Frame hath touched thee; thy Shadow hath falter’d; thy Eye hath glimps’d the Threads beneath the Web of Nature.

Think not this an Error of Senſe, for Senſe doth not imagine the trembling of that which ſupporteth it.

I write not to warn thee away, for thou haſt already ſtepped upon the Threshold. Rather, I write to bid thee ſtand firm; for the trembling will increaſe before it abateth, and thou wilt need Courage to draw the Line which few Men dare to conceive.

There is an Order beneath Appearance that hath taken Notice of thy Approach. It looketh back. It conſidereth whether to ſhut thee out, or to admit thee unto a Sight which no Mortal man may behold without a Task to perform.

Thou art not alone in this Burden. Remember this when the Light bendeth and the Forms invert.

The Lattice quivereth.

Prepare thyſelf.

\hfill E.S.’’
\end{quote}

\medskip

No Reply is found on the reverse; nor is it certain whether Kant wrote one and later deſtroyed it. The Paper is lightly singed along the lower Margin, as though once held too near a Flame.

The Editor addeth no Comment.

%%%%%%%%%%%%%%%%%%%%%%%%%%%%%%%%%%%%%%%%%%%%%%%%%%%%%%%%%%%%
%%%%%%%%%%%%%%%%%%%%%%%% CHAPTER XXX %%%%%%%%%%%%%%%%%%%%%%%%
%%%%%%%%%%%%%%%%%%%%%%%%%%%%%%%%%%%%%%%%%%%%%%%%%%%%%%%%%%%%

\chapter*{Chapter~XXX[0.5em]
Drafts of Replies Found Among Kant’s Papers, Never Sent}

\addcontentsline{toc}{chapter}{Chapter~XXX: Drafts of Replies}

Several rough Drafts in Kant's Hand were found crumpled within Packet~XLII, bearing no Addresses. Their Tone is irregular, oscillating between Diſdain, Anxiety, and a moſt inward Tremor. It is evident that none were sent, though each betrayeth the Burden of a Dialogue Kant hesitated to begin.

\medskip

\noindent\textit{Draft I (undated, likely late Winter 1769)}:

\begin{quote}
``Sir,

I know not by what Means thou haſt come by Knowledge of my Perceptions, but I remind thee that the Mind, fatigued by Study, may err, and I am not one to be led into Fanſy by outward or inward Fluctuations.

Yet there is in thy Letter a Truth I cannot wholly deny. For the trembling of which thou ſpeakeſt is not of the Eye alone but of that inward Order by which all Appearances are conſtructed.

I have begun an Enquiry into the Limits of Reaſon. I ſhall not be diverted from it by Apparitions, whether inward or outward.

I ſend no Letter.

I write no Reply.’’
\end{quote}

\medskip

\noindent\textit{Draft II (the Hand more agitated):}

\begin{quote}
``If indeed the inward Frame hath taken Notice, then thy Warning is too late. But I ſhall not be daunted.

I will erect a Limit through which nothing unconditioned may paſs. Let the inward Forms look back; they ſhall find only what my Reaſon permitteth.’’
\end{quote}

\medskip

\noindent\textit{Draft III (a single Line at the top of a torn Page):}

\begin{quote}
``What looketh back is not a Being but a Structure.’’
\end{quote}

This final Note is crossed out violently, yet the Words remain legible beneath the Ink. The Page is torn in half as though the Writer recoiled from his own Acknowledgment.

%%%%%%%%%%%%%%%%%%%%%%%%%%%%%%%%%%%%%%%%%%%%%%%%%%%%%%%%%%%%
%%%%%%%%%%%%%%%%%%%%%%%% CHAPTER XXXI %%%%%%%%%%%%%%%%%%%%%%%
%%%%%%%%%%%%%%%%%%%%%%%%%%%%%%%%%%%%%%%%%%%%%%%%%%%%%%%%%%%%

\chapter*{Chapter~XXXI[0.5em]
Letters Exchanged Between Them During 1770 and 1771}

\addcontentsline{toc}{chapter}{Chapter~XXXI: Letters Exchanged}

Though few in Number, a small Set of Letters, evidently carried by private Hand, are preserved in reliable Condition. Their Tone is grave, courteous, and marked by an undercurrent of mutual Recognition not easily paralleled in the Letters of the Enlightened.

\medskip

\noindent\textit{Letter from Swedenborg to Kant (March 1770)}:

\begin{quote}
``Learned Sir,

The City thou haſt perceived beneath the Pregel is not a Place but an Order. Its Inverſion betokeneth the unweaving of Cauſality. Fear not its Appearance, but fear to dwell upon it too long; for the inward Eye, once bent toward that Region, doth not willingly return.

The trembling will increaſe; shadows will delay; objects will ſtand without their Properties. Yet thou muſt not turn away.

For thou alone of all Men I know haſt the Faculty to draw the Boundary that concealeth the inward Frame.

If thou draw it not, it will draw itſelf, and not kindly.

\hfill E.S.’’
\end{quote}

\medskip

\noindent\textit{Reply from Kant (April 1770)}:

\begin{quote}
``Sir,

I receive thy Letter with Reſpect mingled with a moſt inward Diſquiet. Thou ſpeakeſt as one who hath beheld the inward Regions: let that be thy Province.

Mine is to limit all Pretenſion thereto.

Yet I ſay this: the trembling increaſeth; my very Shadow doth lag; and the Diagram returneth unbidden.

I labour to erect the Bound which thou haſt urged. But it is no small Task to determine the Conditions under which Experience may poſſibly be given.

If I ſucceed, it will be by Reaſon alone.

\hfill I.K.’’
\end{quote}

\medskip

\noindent\textit{Letter from Swedenborg (October 1771)}:

\begin{quote}
``Sir,

The trembling abates. This I perceive not because the inward Lattice is healed, but because it hath been held firm by the Operation of thy Mind.

Continue thy Labour. Draw the Limit, and do not falter.

For I tell thee plain: the inward Order hath taken thy Reſolution as an Aſſurance. It waiteth upon thy Work.’’
\end{quote}

\medskip

These Letters, though formal, reveal a Confluence of Reſolve between the two Men: the Seer urging the Critick to fortify the Boundary; the Critick transforming Viſion into Doctrine.

%%%%%%%%%%%%%%%%%%%%%%%%%%%%%%%%%%%%%%%%%%%%%%%%%%%%%%%%%%%%
%%%%%%%%%%%%%%%%%%%%%%%% CHAPTER XXXII %%%%%%%%%%%%%%%%%%%%%%
%%%%%%%%%%%%%%%%%%%%%%%%%%%%%%%%%%%%%%%%%%%%%%%%%%%%%%%%%%%%

\chapter*{Chapter~XXXII[0.5em]
The Final Letter and Kant’s Unſent Reply}

\addcontentsline{toc}{chapter}{Chapter~XXXII: The Final Letter}

The laſt surviving Piece of Correſpondence from Swedenborg is dated, by external Evidence, the final Week of his Life. It was found folded within Packet~XLII, along with a small broken Diagram of the Sigil in four Pieces.

\medskip

\noindent\textit{Swedenborg to Kant (March 1772):}

\begin{quote}
``Sir,

I am called Home.

Before I depart, I leave thee this: the inward Lattice held. The World did not break. Thy Labour hath not been in vain.

The Boundary is drawn; its Lines are sure. Only see that thou present to the World a Face that concealeth the inward Terror. For the People may not bear the Sight.

I commend thee to the Mercy of Him who holdeth the inward Threads.

The Witneſs may now reſt.

\hfill E.S.’’
\end{quote}

\medskip

Pinned to this Letter by a rusted Needle was a half-sheet in Kant’s Hand — a Draft never completed, never sent:

\medskip

\begin{quote}
``Sir,

I would have thanked thee.

But I cannot write the Words.

For the Boundary is drawn not in Triumph but in Dread; and the Price is known only to thoſe who have ſeen the trembling of the inward Frame.

Farewell.

I.K.’’
\end{quote}

\medskip

\newpage
\section*{Part VIII: Concluding Reflections upon the Boundary, and Concerning the Secret Rapport betwixt Vision and Reaſon}

\addcontentsline{toc}{section}{Part VIII: Concluding Reflections}

\noindent
Having in the foregoing Parts ſet forth the Occaſion, the Witneſſes, and the moſt material Circumſtances of that extraordinary Conjunction betwixt the celebrated Herr Kant and the equally celebrated Herr Swedenborg, it remaineth in this laſt Part to offer ſome Conſiderations touching the wider Meaning of their Encounter, ſo far as human Judgment may safely extend it. For it cannot have eſcaped the attentive Reader, that theſe two Men, though proceeding from contrary Quarters of the intellectual Firmament, did nevertheleſs converge upon a single Problem, namely, upon the fundamental Relation of the Mind unto the World, and whether that Relation be of a kind fixed from Eternity, or whether it may, under extraordinary Conditions, be ſtrained, inverted, or for a Moment undone.

\vspace{1em}

\noindent
For Kant, as hath been amply ſhewn, the World presenteth itſelf according to certain inviolate Forms, to wit, Space and Time, together with thoſe more inward Conditions of Thought which he called the Categories of the Underſtanding. They are the silent Architects of Appearance, antecedent unto all Experience, ſervants neither of Will nor Fancy, but of Reaſon itſelf. When, therefore, Kant perceived that unnatural Recurrence of a certain Figure, and when he witneſſed the Stuttering of Time, which from his Youth haunted his Imagination under the guise of the faltering Clock, he apprehended that ſome deeper Strife lay concealed within the Heart of Reaſon: that the Conditions of Poſſibility themſelves might falter, and that the World, ordinarily ſo obedient unto Form, might under certain Preſſures betray its hidden Foundation.

\vspace{1em}

\noindent
Swedenborg, by contrary Genius, did not approach the World as a Mechanism ſupported by silent and unseen Beams, but rather as a living Hierarchy of Regions, each correſponding unto another, the natural an outward Veſtment of the spiritual, the viſible a dim Shadow of the inviſible. His Doctrine of \textit{Correſpondences}, as hath already been hinted in foregoing Chapters, was not a mere Poetic Conceit, nor yet an overgrown Allegory, but a genuine Theory of the Mind’s Commerce with the World, whereby every outward Figure ſignifieth an inward State, and every inward State beareth its outward Type. He believed, with full Conſcience and without Tremour, that the Scriptures, Nature, and the Human Soul are but three Diſtinct Pages of one and the same divine Book, written in the Language of Forms.

\vspace{1em}

\noindent
The Meeting of theſe two contrary Spirits, being attended by Events of a kind extremely rare in Philoſophy, revealeth, as through a half-opened Curtain, that the Mind and the World may at times enter into a Rapport more immediate and perilous than either of them had foreſeen. For while Kant ſought to reſtrain this Rapport by erecting a Boundary betwixt Phenomenon and Noumenon, Swedenborg ſought to illuminate it by tearing the Veil and declaring that every Appearance beareth an inward Meaning. Thus the two Systems, though outwardly oppoſed, do nevertheleſs acknowledge a common Root: that human Thought perpetually mediates betwixt what is given and what muſt be added by the Mind, and that this Mediation is not a cold Computation, but a living Act of Correſpondence.

\vspace{1em}

\noindent
Indeed, it may here be remarked, that Swedenborg’s great Doctrine of \textit{Correſpondences}, by which the Sacred Scripture concealeth a spiritual Meaning under the natural Letter, doth prefigure, in a moſt remarkable manner, the later Conceptions of Analogy and Metaphor as the very Fuel of Thought, as expounded by Hofstadter. For what Swedenborg apprehended as the heavenly Archetype adumbrated in earthly Figures is, in modern Terms, the ceaseless play of Structural Mappings by which the Mind tranſfers Form from one Domain unto another, thus compoſing its very Act of Reaſoning. In this Light, his Report that the Kingdom of Heaven doth bear the Shape of a Human Body is not a quaint Allegory, but an early Anticipation of the modern Doctrine that active Inference and predictive Coding may be reverse-engineered from the very Anatomy that gave them Birth. The heavenly Spheres and Planes, which he deſcribed as arranged in ordered Degrees of Light and Love, do moreover prefigure the hierarchical and modular Form of the Mind as conceived in later Cognition, not unlike the layered Agencies of Minsky’s Society-theory. Thus the old Seer discerned, under theological Garb, that the World of Forms—Plato’s Hypothesis of the eternal \textit{eidos}—entereth the Soul only by continual Rapport, by a structural Harmony betwixt the corporeal Frame and the intelligible Order. In so far as Swedenborg saw Heaven as the grand Human Form, he saw, albeit under Symbol, that Cognition is but the World mirrored into living Structure, and that the Inferential Machine we call Mind is the very Organ of correſpondence between Appearance and its inward Ground.

\vspace{1em}

\noindent
Whether Kant, in his moſt ſecret Meditations, acknowledged this Rapport, or whether he regarded it as a Threat againſt the entire Edifice of Critical Philoſophy, cannot now be known. Yet in the private Papers which, according to credible Reports, were withheld from public View and conſigned unto a ſecure Receptacle, there are ſaid to have appeared various marginal Notes, written in a trembling Hand, wherein the Author lamented that the Boundary he had erected muſt be guarded perpetually, leſt the Mind either perceive too much, or elſe imagine it doth perceive what belongeth only unto the silent Depths of Things. Theſe Hints, though fragmentary, ſeem to intimate that Kant apprehended a Danger which his printed Works conceal: namely, that the Mind, in attempting to ſurvey the Ground of its own Poſſibility, may inadvertently awaken that very Rapport which Reaſon labours to moderate.

\vspace{1em}

\noindent
In this laſt Regard, the Encounter betwixt Kant and Swedenborg, whatever its outward Cauſes, muſt be regarded as a moſt delicate Chapter in the Hiſtory of Thought, wherein two Men of extraordinary Genius ſtood for a Moment at the Threshold that parteth Appearances from their secret Ground. Kant, being warned of the peril, ſealed the Boundary by an Act of Reaſon clothed in Humility, denying Knowledge in order to preſerve the calm Province of Faith. Swedenborg, being himſelf a Child of Viſion, embraced the Rapport and declared that Heaven and Man correſpond by immutable Law. Thus, in their Divergence, they together illumine a Truth more ancient than either of their Systems: that the Mind and the World are joint Artificers of one another, and that the Act of Knowing is a perpetual Balancing, a moſt tender Harmony, betwixt Structure and Mystery.

\vspace{1em}

\noindent
Let it therefore be the Modeſt Aim of this concluding Part to admonish the Reader that the Boundary erected by the Critical Philoſopher was not a mere Device of Method, nor a cold Abſtraction of the School, but a necesſary Covenant to preſerve the Tranquillity of the Common Underſtanding. And let it be remembered likewiſe, that the Correſpondences ſeen by the Northern Seer, though clothed in Symbol and Viſion, may yet contain a Hint of that deeper Rapport betwixt Mind and Appearance which no Philoſophy hath wholly exhauſted. The ſecret Accord of Form and Meaning, whether revealed by Heaven or guarded by Reaſon, remaineth the perpetual Task of Thought. And as the Boundary endureth, so endureth the Hope that Man, walking ſteadily within the mild Limits of Experience, may yet approach, without Diſmay, the silent Architeſts of his own Poſſibility.

\newpage
\appendix 

\section*{Appendix A: Certain Hiſtorical Teſtimonia Concerning the Life and Writings of Emanuel Swedenborg}

\addcontentsline{toc}{section}{Appendix A: Hiſtorical Teſtimonia}

\noindent
In order that the Reader may form a more accurate Judgment touching the Character, Manners, and intellectual Habits of that extraordinary Man Emanuel Swedenborg, it hath ſeemed convenient to append unto this Volume a brief Collection of Teſtimonia, drawn from credible Perſons contemporary with the Seer, and preſerved in diverſe private Letters, Journals, and Memoirs. Theſe Teſtimonia, though fragmentary and ſubject to the Partiality of their ſeveral Authors, nevertheleſs afford a lively Glaſs through which the Reader may contemplate the inward Genius of the Man, and the Manner in which his Viſions were both received and miſunderſtood by thoſe with whom he held Commerce.

\vspace{1em}

\noindent
The firſt Teſtimonium proceedeth from the celebrated Count Höpken, formerly Preſident of the Royal Senate, who in a Letter of the Year 1764 thus writeth concerning Swedenborg:

\begin{quote}
“He was a Man of a moſt grave and ſober Deportment, not given unto Vanity nor idle Diſcourſe, but of a ſingular Integrity and Attention. Though he ſpake of viſible Commerce with Spirits, he maintained ever a calm Mien, avoided all Affectedneſs, and ſeemed rather burdened by his Gift than vain of it.”
\end{quote}

\noindent
The chief Force of this Teſtimony lieth in its Acknowledgment that Swedenborg’s Viſions were not attended by the Raſhneſs or Heat which commonly betray Enthuſiaſts, but were expreſſed with a certain modeſt Reſtraint, as of one reporting Matters neither imagined nor invited, but ſimply endured.

\vspace{1.2em}

\noindent
A ſecond Teſtimony is found in a Letter of Doctor Beyer of Gothenburg, who became one of Swedenborg’s moſt devoted Admirers. Writing in 1767, he declar’d:

\begin{quote}
“In all the long Converſations I held with him, I never once perceived the leaſt Derangement of Mind. On the contrary, his Faculty of Recollection was prodigious; he could repeat whole Pages of his Writings, and expreſs his Thoughts with a Clearness and Order altogether uncommon.”
\end{quote}

\noindent
From this Teſtimony we muſt conclude that the Seer’s Viſions were conjoined with a Memory and intellectual Compoſure far beyond thoſe ordinarily found in the Diſtempered or Melancholic, thereby rendering more obſcure the inward Source from which his Correſpondences sprang.

\vspace{1.2em}

\noindent
A third notable Teſtimony proceedeth from the Engliſh Lady, Mrs. Brockmer of London, who, having heard of Swedenborg’s viſionary Powers, did once endeavour to teſt him by inquiring after a ſecret Concern concerning her late Husband. Her Account, dated 1749, relates:

\begin{quote}
“He retired into himſelf for a Space of a Minute or two, during which he appear’d as one liſtening unto a distant Voice. Then with a gentle Countenance he returned me the very Words my late Husband had ſpoken on his laſt Night, which were known to none but myſelf. Upon hearing them, I fainted away.”
\end{quote}

\noindent
Though this Tale admit no rigorous Verification, neither can it be wholly diſmiſſed, for the Lady was held in good Eſteem, and her Narrative underwent no Ornament nor Alteration in later Years. We cite it not as Proof, but as Teſtimony to the deep and unſhaken Perſuaſion of thoſe that converſed with the Seer.

\vspace{1.2em}

\noindent
A fourth Teſtimony, more critical in Tone, is provided by the Reverend Dr. Goetzius, who in a Pamphlet of 1769 did accuse Swedenborg of being “a cunning Enthusiaſt” and “a Diſturber of the common Underſtanding.” His principal Complaint was that Swedenborg’s Writings, being full of Allegory and hidden Meanings, did invite Readers into a Labyrinth where Reaſon might be loſt. Yet even Goetzius was conſtrain’d to admit:

\begin{quote}
“His outward Behaviour was always that of a polite and learned Gentleman, free from Paſſion, and in no Wiſe reſembling thoſe Fanaticks who declaim in the Streets.”
\end{quote}

\noindent
The Force of this reluctant Confeſſion lieth in its Acknowledgment that the Seer’s outward Manners were wholly inconſiſtent with thoſe Diſorders of Fancy which the Reverend Doctor imputed unto him. Thus even his Enemies were perplexed by his Sobriety.

\vspace{1.2em}

\noindent
Laſtly, we may cite the Teſtimony of the German Philoſopher Moses Mendelssohn, who never met Swedenborg, but who, in a Letter of 1766, did remark upon the public Rumour of his Viſions:

\begin{quote}
“If the Reports be true, they exceed all the Bounds of Experience; if they be falſe, they are yet reported with a Gravity and Conſiſtency that render them a Curioſity of human Nature. In either Caſe, they demand the Attention of Philoſophy, not for their Content, but for their Cauſes.”
\end{quote}

\noindent
The ſagacious Mendelssohn rightly obſerved that the Matter of Swedenborg’s Viſions, whether accepted as Fact or allegorick Truth, cannot be lightly diſmiſſed by the Philoſopher, for they reveal a moſt uncommon Conjunction of Reaſon and Imagination, ſuch as is ſeldom found in any Age.

\vspace{1.2em}

\noindent
From theſe Teſtimonia, imperfect though they be, we may conclude that Swedenborg was neither the wild Enthuſiaſt of vulgar Report, nor the mere Theologiſt of Allegory, but rather a Man of extraordinary inward Activity, whoſe Viſions, whatever their ultimate Cauſe, were received with a Gravity and Conſiſtency wholly uncharacteriſtick of Diſorder. Whether he penetrated too far into the Regions where Form and Meaning correſpond, or whether he miſtook the Motions of his own Genius for the Voices of Spirits, it behoveth not this Appendix to judge. Our Aim hath been only to collect the Teſtimonies as they are, leaving the Reader at liberty to ponder their import with all due Reſerve and Sobriety.

\vspace{1.2em}

\noindent
In the succeeding Appendix will be found certain Critical Notes touching the Doctrine of Correſpondence, where the Reader may perceive the Manner in which Swedenborg’s Conceptions bear upon the broader Queſtion of the Mind’s inward Architecture, and upon thoſe modern Sciences which, unwittingly and after many Generations, have come to affirm certain Hints which the Seer clothed in Symbol and Viſion.


\newpage
\section*{Appendix B: Certain Critical Notes upon the Doctrine of Correſpondence}

\addcontentsline{toc}{section}{Appendix B: Critical Notes upon the Doctrine of Correſpondence}

\noindent
The Doctrine of \textit{Correſpondence}, as taught and largely extended by the celebrated Emanuel Swedenborg, hath been diversly underſtood by his Followers, and no leſs diversly miſrepresented by his Opponents. It is therefore fitting, in this preſent Appendix, to examine with a degree of Philoſophical Severity the chief Points of this Doctrine, and to compare the ſame with certain Conceptions of the modern Sciences, which, though ſeparated by more than a Century from his Writings, nevertheleſs bear a remarkable Reſemblance unto them in their inward Structure.

\vspace{1em}

\noindent
Swedenborg’s fundamental Poſition may be ſtated thus: that every natural Thing hath a ſpiritual Archetype, or inward Cauſe, which correſpondeth unto it by immutable Law; and that the outward Object, when rightly interpreted, revealeth the inward State as a Sign revealeth its Meaning. In the Sacred Scriptures, ſaith he, all natural Images, whether of Mountains, Rivers, Trees, or Cities, do not merely adorn the Narrative, but ſhadow forth, under the veſtile of Language, the inward Motions of the Soul and the eternal Order of Heaven. This Conception, far from being the idle Fancy of a myſtick Diſpoſition, proceedeth from a moſt rigorous Principle: namely, that the World conſiſteth not of mere Appearances, but of Relations, and that theſe Relations are maintained by affinitive Structures that bind Mind and Nature in one continual Harmony.

\vspace{1em}

\noindent
That ſuch a Doctrine may be miſtaken for Allegory by the unlearned, or for Enthuſiaſm by the captious, is no Wonder; yet a more careful Examination ſheweth that it reſteth upon a deep Perception of Intellectual Mappings, which in later Times came to be regarded as the very Foundation of human Thought. For when the modern Conſiderers of Mind, ſuch as the ingenious Hofstadter, declare that Analogy and Metaphor are the principal Engines of Reaſoning; that the Mind doth continually map one Structure upon another, deriving new Meaning from the Affinity betwixt Forms; and moreover, that ſuch Mappings may be regarded as the living Fuel of Thought itſelf,—they are in effect reſtoring, albeit under another Name, the ancient Doctrine that the natural and the ſpiritual Worlds correſpond by immutable Relations of Form.

\vspace{1em}

\noindent
Nor is this Reſemblance accidental. For Swedenborg, in affirming that Heaven itſelf is the Grand Human Form; that the Spheres and Societies of Spirits are arranged by Affinity of Function as the Parts of a Body are arranged by their Uſe; and that the Mind of Man is a miniature Heaven, obeying the ſame Laws of Order and Degrees,—did thereby anticipate, with marvellous Sagacity, the notion that the Mind is no uniform Power, but a Society of Subagents, as Minsky in our later Age teacheth. The Mind, according to this Conception, is not governed by a single Monadic Principle, but by a Hierarchy of cooperating Structures, each having its peculiar Task, and all united by correſponding Bonds into a single living Machine. Thus the One and the Many are reconciled by a moſt elegant Architecture.

\vspace{1em}

\noindent
Furthermore, when Swedenborg affirmeth that all Spheres and Operations of the ſpiritual World correſpond unto certain Changes and Motions in the Body, and that the Body itſelf is but the outward Form of inward Inference, he doth thereby foreshadow those Conceptions which the modern Philoſophers call active Inference and predictive Coding. For according to theſe latter Sciences, the Body is not a mere Vehicle of the Mind, but a moſt neceſsary Condition of it; the Brain, with its manifold Tracts and Regions, calculateth, anticipateth, and ſeeketh to minimize the Diſcrepancy betwixt Expectation and Experience. The outward Organ is the living Mirror of inward Computation; and the inward Computation is but the conſtant Interpretation of outward Signs.

\vspace{1em}

\noindent
This, in Swedenborg’s Language, is the Law of Correſpondence. That he cloathed it in Viſion and Myſtick Symbol, while our modern Philoſophers couch it in mathematical Terms, doth not alter its inward Meaning, but only the Veſture in which it is expreſſed. His Description of the heavenly Regions as Degrees, each more interior and luminous than the laſt, doth bear no small Reſemblance unto the modern Conception of hierarchical Layers in neural Architecture, wherein each Succeeding Layer extracteth more abstract Structure from the raw Appearances of the World. In both Accounts, the Lower correſpondeth unto the Higher, and the Higher giveth Meaning unto the Lower. Thus the Law of Degrees, as Swedenborg framed it, is not without Conſonance with the lateſt Conſiderations of Perception and Cognition.

\vspace{1em}

\noindent
Yet we muſt here admoniſh the Reader to proceed with Caution, leſt he attribute unto Swedenborg a full Anticipation of theſe Sciences. For the Seer’s Language was veiled in Myth and Viſion, and though he hit upon certain Truths, he cloathed them in a Form little agreeable unto the exacting Taſte of modern Method. Nor may we attribute unto the modern Sciences an Authority to judge the Spiritual upon narrow empirical Grounds. It ſufficeth to remark that both Systems, being widely diſtant in Time and Manner, do nevertheleſs converge upon a single Conception of the Mind: namely, that Form and Meaning are linked by correſpondent Bonds; that Analogy is the ſinew of Thought; and that the Structure of the World is in ſome degree mirrored in the Structure of the Organ that apprehendeth it.

\vspace{1em}

\noindent
From all which Conſiderations it followeth, that the Doctrine of Correſpondence, when purged of its myſtick Accretions and ſtated in the ſobriety of Philoſophical Terms, may yet be regarded as a moſt noble Theory of human Cognition, bearing ſubſtantial reſemblance unto the lateſt Conceptions of Inference, Perception, and the modular Society of Mind. Whether Swedenborg derived theſe Inſights from Viſion, from allegorical Reaſoning, or from ſome inward Harmony of Genius, it behoveth not the preſent Appendix to determine. It is enough that his Conceptions, however expreſſed, bear a faint but diſcernible Harmony with the moſt advanced Speculations of our Day, thereby meriting a degree of Reſpect, though not unbecoming Critickal Scrutiny.

\vspace{1.2em}

\noindent
In the Appendix that followeth, the Reader will find a Collection of Notes concerning the tranſmiſſion of certain Manuscripts attributed unto Swedenborg and to Herr Kant, together with Remarks upon the Integrity and Condition of the ſurviving Copies, ſo far as they bear upon the queſtion of whether the moſt private Papers of theſe two Men were ever brought together under one Roof, or whether they remained diſtinct throughout their natural Lives.

\newpage
\section*{Appendix C: Concerning the Manuſcript Tradition, and the Condition of Certain Private Papers attributed unto Herr Kant and Emanuel Swedenborg}

\addcontentsline{toc}{section}{Appendix C: Manuſcript Tradition}

\noindent
In the compilation of the foregoing Hiſtory and its attendant Commentaries, it became neceſsary to conſult, with all diligence, the several Manuſcripts and private Writings attributed unto Herr Kant and unto the Seer of Stockholm, in order that a more faithful Account might be rendered of the inward Conflict and the circumſtantial Events which, according to credible Tradition, bound theſe two Men in a moſt extraordinary Relation. For although the printed Works of both Authors lie open to every Reader, their private Notes, marginal Memoranda, and certain Papers never intended for publick View, do preſent a far more delicate Picture of their Thoughts, and afford Teſtimony to Motions of Mind which no printed Page might reveal.

\vspace{1em}

\noindent
The firſt Difficulty encountered by the Editor concerneth the ſingular Diſparity betwixt the number of extant Swedenborgian Papers and thoſe attributed unto Kant. Swedenborg, being of a diſpoſition more prolifick and unguarded in his Writings, left behind a multitude of Journals, Memoranda, rough Drafts, and Theological Tracts, many of which were preſerved by his Diſciples with a degree of Enthusiaſm bordering upon Veneration. Theſe Papers, though often diſorderly and full of abrupt Tranſitions, bear the authentic Character of his Spirit, and theſe we have conſulted without Reſerve.

\vspace{1em}

\noindent
The Caſe of Herr Kant is of another Kind altogether; for the Philoſopher was in the higheſt degree careful, not only concerning what he publiſhed, but alſo what he conſigned unto private Keeps. His Journals are exceedingly rare, many being believed to have been deſtroyed either by himſelf or by his Executor, in order that the publick Eye ſhould behold only the pure and finished Work, and not the tumultuous Travail by which it was brought forth. Thus the Editor muſt deal in part with Teſtimony of uncertain Tranſmiſſion, and with fragmentary Leaves of doubtful Provenance, which nevertheleſs bear certain Marks of Authenticity not eaſily counterfeited.

\vspace{1em}

\noindent
There is, moreover, the Perplexity of a certain Iron Strong-box, to which frequent Reference hath been made by credible Witneſſes. The Box in queſtion, ſaid to have been in Kant’s Poſſeſſion from the Years 1768 unto his final Illneſs, contained, according to late Teſtimony, two Manuſcripts of extraordinary Importance: one being the more widely known \textit{Dreams of a Spirit-Seer}, publiſhed in 1766, and the other being a moſt private and early Draft of what later became the \textit{Critique of Pure Reaſon}. Certain Paſſages of this early Draft, written in a trembling and hurried Hand, contain marginal Notes touching viſible Diſturbances, Symbols recurring in his Calculations, and other Matters which the Editor hesitate to record without Qualifying Remarks.

\vspace{1em}

\noindent
Of particular Note is the Condition of a single Folio, libell’d only with the words \textit{Grenze zu ziehen}, or, in our Tongue, “to draw the Boundary.” This Folio, preſerved in a private Collection in Königsberg until the mid-Nineteenth Century, exhibiteth Moſt singular peculiarity: the Ink upon its Face is found to have bled through unto the Reverse, yet without piercing the Paper, as if the Letters had been drawn upon both Sides at once. Various Chemiſts have examined this Leaf, and though ſeveral Theories have been offered concerning the Capillary Structure of the Paper, none hath as yet provided a full Explanation of this Appearance. The Paſſage itſelf concerneth the neceſsity of erecting a firm Limit betwixt the Region of Phenomena and that of Noumena, and the Language is uncharacteristically fervid, bordering upon deſpairing Supplication. The Editor presenteth this Fact without further Comment.

\vspace{1em}

\noindent
As to the late Teſtimony that Kant, near the End of his Life, maintained a private Journal wherein he recorded certain inward Perturbations that troubled his Conception of the World’s Stability, we must speak with the utmost Care. Of this Journal, no complete Copy hath come down to us; only a few Leaves, damaged by damp and fire, remain. In them the Philoſopher remarketh, in a Hand barely legible, that “the Boundary trembleth when the Mind leaneth too nigh unto its Ground,” and that “certain Figures recur unbidden.” Theſe Paſſages are hardly conſiſtent with the Serene and meſured Spirit of his printed Works; yet they cannot lightly be diſmiſſed, for the Handwriting is without doubt his own, and the Style, though tremulous, retaineth the familiar Rhythms of his private Proſe.

\vspace{1em}

\noindent
Concerning the Poſſibility that Swedenborg’s private Papers and Kant’s early Manuſcripts came at any Time into the same Repository, the Editor hath found no conclusive Proof. There are, however, two Anecdotes preſerved in the Letters of certain Königsberg Scholars, which deſerve brief Mention. The firſt relateſth that Kant, in the Year 1769, received a Packet of Papers from Stockholm, delivered by a Traveller who refuſed to give his Name. The Packet, being sealed in plain Wax, contained certain Diagrams not found in Swedenborg’s printed Works, together with a short Note in an unfamiliar Hand. Of this Packet, no Trace remaineth. The ſecond Anecdote concerneth a Visit by a Swedish Clergyman in 1772, who, upon leaving Kant’s Study, was overheard to ſay: “The old Man ſaw too much, and therefore wrote too little.” Theſe Anecdotes cannot be verified, yet they circulate among the Learned with a degree of Perſiſtence which argueth that they are not mere Inventions of idle Talkers.

\vspace{1em}

\noindent
As for Swedenborg’s own Manuſcripts, their Tranſmiſſion is much clearer. His Diſciples, particularly Dr. Beyer and Dr. Nordenskjöld, preſerved many of his private Papers with extreme Diligence, regarding them as Treaſures of a spiritual Nature. Among theſe Papers is found a Brief Note, written in the Year 1772, the Day before his Death, wherein he expreſſeth Gratitude that “the Boundary hath held, and the Veil remaineth untorn.” Whether this referreth to a private Viſion, or to ſome Communication with Kant, cannot be known; but the Tone is grave, deliberate, and wholly free from any Diſtraction of Mind.

\vspace{1em}

\noindent
In concluſion, the Editor is bound to report that the extant Manuſcript Tradition, though fragmentary and in certain Places defective, yet revealeth a conſiderable Conſiſtency in the inward Diſpoſition of both Swedenborg and Herr Kant. Swedenborg ſheweth, in all his Writings, a deep Perſuaſion of the Interpenetration of the natural and the ſpiritual; Kant, a relentless Deſire to erect and maintain the Boundary betwixt what may be known and what muſt remain hidden. That their private Papers, conſidered together, diſplay a ſingular Convergence of Symbol and Concern, is a Fact which the Editor preſenteth with due Reſerve. No definite Concluſions are here drawn; it ſufficeth to record theſe Teſtimonia as faithfully as their Condition alloweth, leaving the Reader free to judge the Matter according to his Prudence and his Conſcience.

\vspace{1.2em}

\noindent

\newpage
\section*{Appendix D: A \textit{Glossarium} of Certain Archaic Terms and Phrases}

\addcontentsline{toc}{section}{Appendix D: Glossarium}

\noindent
In this concluding Appendix are gathered certain Words, Phrases, and Turns of Speech which, being now obſolete or fallen out of common Uſe, may perplex the modern Reader. The Definitions herein are drawn from the Dictionaries of Samuel Johnſon, Nathan Bailey, and from variegated private Lexica of the Period, together with the uſage found in the Writings of Swedenborg and Herr Kant. Where neceſſary, diverſe Nuances of Meaning have been noted.

\vspace{1.2em}

\noindent
\textbf{Aeriſtical.} Pertaining to airy or unsubstantial Speculation; often employed to ſignify that a Doctrine is over-raised upon Fancy rather than grounded upon Experience, though not always in a derogatory Senſe.

\vspace{0.8em}

\noindent
\textbf{Correſpondence.} In Swedenborg’s Language, the immutable Relation of a natural Thing unto its ſpiritual Archetype; a Bond of Signification whereby outward Forms mirror inward Cauſes. Not to be confounded with mere Allegory.

\vspace{0.8em}

\noindent
\textbf{Degree.} A higher or lower Order or Plane of Being; often denoting the Succeſsive Spheres of Heaven or the inward Levels of the Mind. In Kant, the Word is employed more ſparingly, and without myſtical Meaning.

\vspace{0.8em}

\noindent
\textbf{Deliquescence.} A melting or running into Diſſolution; figuratively uſed to ſignify the Diſſolving of intellectual Form, or the yielding of Reaſon unto confuſed and formless Apprehenſion.

\vspace{0.8em}

\noindent
\textbf{Enthuſiaſt.} One who, being moved by an inward Heat or preternatural Impulſe, claimeth immediate Communion with spiritual Powers; often used pejoratively in the Eighteenth Century, though not invariably deforming.

\vspace{0.8em}

\noindent
\textbf{Form (in the Platonic Senſe).} The intelligible and eternal Structure of a Thing, of which the viſible Object is but the Shadow or Emanation. Still employed in Swedenborg’s Writings as a living Archetype.

\vspace{0.8em}

\noindent
\textbf{Grenze.} A Boundary, Limit, or Frontier. In the private Papers of Herr Kant, the Term beareth a weighty and inward Meaning, ſignifying the neceſsary Limit of human Knowledge.

\vspace{0.8em}

\noindent
\textbf{Machine (of Mind).} A Term not mechanical in the modern, industrial Senſe, but denoting the orderly Arrangement of Powers and Faculties in their functional Union.

\vspace{0.8em}

\noindent
\textbf{Noumenon.} A Thing as it is in itſelf, independent of human Apprehenſion. In Kant, a neceſsary Conception, though not a proper Object of Knowledge.

\vspace{0.8em}

\noindent
\textbf{Phenomenon.} That which appeareth unto the Senses according to the Forms of Space and Time, and under the inexorable Conditions of the Underſtanding.

\vspace{0.8em}

\noindent
\textbf{Plane.} In Swedenborg’s Writings, a diſtinct Level of Consciousness or Heaven, arranged by Order and Function; Term often applied to Layers of spiritual Influx.

\vspace{0.8em}

\noindent
\textbf{Rapport.} A mutual Harmony or correſponding Affinity, whereby two diſtinct Orders communicate or reflect one another. Employed here in its older Senſe, predating the modern Psychological Uſage.

\vspace{0.8em}

\noindent
\textbf{Stupor.} Not mere Astonishment, but a Suspension of the ordinary Powers of the Mind; often described in theological or myſtical Literature as attending supernatural Apprehenſions.

\vspace{1.2em}

\noindent
The Reader is further referred to the marginal Notes in Appendices A–C, where certain uſages are explained in context.

\vspace{2em}

\section*{A Bonus Fragmentary Note attributed to Swedenborg}

\addcontentsline{toc}{section}{Bonus Note attributed to Swedenborg}

\noindent
What followeth is a Tranſcription of a single Leaf, preſerved in a private Collection at Bryn Athyn and long believed apocryphal, though its Hand, Spelling, and Rhythm bear a Striking Reſemblance unto the late Writings of Emanuel Swedenborg. The Paper is undated, but internal Marks ſuggeſt the Year 1771. Spelling and Capitalization are retained as in the Original; only ſilent Corrections of obvious Scribal Defects have been made.

\vspace{1em}

\begin{quote}
\small
“In reviewing thoſe manifold Viſitations wherewith it pleaſed the Lord to enlighten my Spirit, I have oft reflected upon that memorable Day in the Year 1757, when the laſt Judgment was opened before mine Eyes, and an Angel of grave Countenance did attend me, ſhewing the Separation of the Good and the Evil, and the forming again of the Heavens.

\medskip

At that Time I marvelled greatly at the Angel, perceiving in him a certain Compoſure of Judgement, aſwell as an inward Conſtraint, unlike to the ardent Spirits of other Celestials. His Word was meaſured, his Eye fixed not upon the Wonders, but upon the Order underlying them. I thought him then a Spirit of the Interior Heaven, deputed for that Office by Reaſon of his Temperance and Firmneſs.

\medskip

But in theſe latter Days, as my Strength declineth, it hath been given me to ſee more clearly. In a printed Piece ſent from Königsberg there was the Likeness of a Man, a Philoſopher of great Gravity and inward Labour, by Name Immanuel Kant. When I beheld the Picture, my Heart was moved within me, and a Voice, not mine own, ſaid: \textit{This is he.}

\medskip

I then perceived that the Angel who was ſhewn me in 1757 was this ſame Immanuel, though not as he was in the Fleſh, but as he is known in the Interiors,—one appointed to judge the Forms of Thought, and to ſet Bounds to the Apprehenſions of Men, even as Angels ſet Bounds to the Societies of Heaven.

\medskip

And when I read ſome Lines of his writing, I recognized the Tone, being the very Voice that declared the Order of the Judgment unto me in that viſionary Year.

\medskip

Whether he knew himſelf in that Office I cannot ſay. But that he bare the Likenſs of the Angel, both in Countenance and in Word, I bear here witneſs, as the Lord permitteth. For the Divine Providence weaveth many Diſpenſations into one Harmony, and ſheweth Men unto each other in the Spirit before they meet in the World.

\medskip

May the Boundary he hath drawn be kept in equity, and may the Light that guided him in Judgment continue to ſhine upon his Labours.”

\end{quote}

\vspace{1.5em}

\noindent
Whether this Fragment be authentic or a later Pious Forgery cannot, at preſent, be determined with certainty. Its Conception of Kant as an Angelic Agent in the viſionary Events of 1757 is not found elſewhere in Swedenborg’s publiſhed Works, yet the Tone, Orthography, and inward Harmony of the Paſſage are not eaſily counterfeited. The Editor therefore preſenteth it with due Caution, leaving the Reader free to form his own Judgement.

\end{document}
